% =====================================================================
%  Harald Høffding, „Positivisme og Materialisme. Med særligt Hensyn til
%  Maudsleys Foredrag ...“, in two parts.
%  Nyt dansk Maanedsskrift, udg. af Vilhelm Møller, 2. Bind (April --
%  September 1871): pp. 193--208 (sections 1--2) and pp. 309--333.
%  TRANSCRIPTION (Danish) -- diplomatic, from the page images.
%
%  ---- SOURCE -----------------------------------------------------------
%  Internet Archive item nyt-dansk-maanedsskrift-bind-2 (whole volume,
%  Nyt_dansk_Maanedsskrift_bind2.pdf, 581 PDF pp., 327118779 bytes, sha256
%  449b659b6c4b4fac3d76ee8835c55f0023536cfb83e6123988ead21b9870cdcb,
%  ABBYY FineReader text layer), downloaded 2026-09-27:
%  https://archive.org/details/nyt-dansk-maanedsskrift-bind-2
%  The two parts are volume PDF pp. 199--214 and 315--339.
%  Kept as two extracts (qpdf, text layer kept) in ~/bibliotek/Høffding, Harald/:
%    1871-positivisme-materialisme-I-nyt-dansk-maanedsskrift-2.pdf   sha256 f8d9dcd5...
%    1871-positivisme-materialisme-II-nyt-dansk-maanedsskrift-2.pdf  sha256 6838c83e...
%  Each page one 300 ppi RGB JPEG.  Fraktur.
%
%  ---- PAGE MAP ---------------------------------------------------------
%  Volume PDF = printed + 6 throughout; verified by eye on the folios.
%
%  ---- CONVENTIONS ------------------------------------------------------
%  Diplomatic; 1871 orthography as printed.  Letterspacing -> \emph;
%  antiqua against the Fraktur -> \textit.  Fraktur I/J resolved by the
%  word.  Quotation marks „...“ as printed.  Footnotes *) at the mark.
%  Printer's errors transcribed as printed and logged in a % comment at
%  the site.
%
%  ---- LAYOUT -----------------------------------------------------------
%  Part 1 opens a fresh page (p. 193, no folio) under the title in display
%  blackletter; the subtitle in heavier Fraktur of body size, with the
%  footnote „Senere udkommet i særligt Aftryk.“; a short rule; „1.“.
%  Section 2 begins p. 200 (no rule).  Part 1 closes (p. 208) with
%  „(Fortsættes.)“, the signature „Harald Høffding.“ set right in bold
%  Fraktur, and a double rule.  Part 2 opens a fresh page (p. 309) under
%  the same title, a rule, John Stuart Mill's English epigraph in small
%  antiqua, a second rule, and „3.“; section 4 begins p. 322 after a rule.
%  It closes (p. 333) with the signature and a double rule.  Each part
%  opens with an enlarged initial.  14 footnotes, some running over onto
%  the next page (given whole at the mark, with a comment).  German
%  quotations are set in Fraktur; French, Latin and English in antiqua
%  (\textit), French quotations within roman guillemets »...« as printed.
%  Emphasis (letterspacing, measured within the line): 32 runs.
%
%  Printer's errors and oddities logged at the site: p. 201, 207 (x3)
%  French quotations closed with » for «; p. 204 „Sammensat“ (capital
%  mid-sentence); p. 206 a closing » with no opening mark; p. 312 no comma
%  after „Materialister“; p. 323 „Geschichte der Materialismus“; p. 327
%  „theoride“; p. 329 a single low ‚ opening „Arvelighed“ (hence the
%  „/“ imbalance of one); p. 314 „ellér“ (a foul-case é); p. 318 a wrong
%  sort for the final e of „tidligere“ (word kept).  Foreign forms as
%  printed: p. 198 „mekanischen“, p. 199 „Rechtsbewußtseins“, p. 201
%  „caracterisé“, „rationelle“, „2éme“.  check.py flags p. 321 „Fornuft-“:
%  a suspended hyphen („Fornuft- eller Bevidstheds-Elementet“), correct.
%
%  ---- ACCURACY  [2026-09-27] ------------------------------------------
%  Word-accuracy: every batch diffed against the embedded ABBYY layer
%  before splicing (ocrdiff.py --frag; offsets -192 and -292 in the
%  41-page working extract): 5 batches (pp. 193-200, 201-208, 309-316,
%  317-324, 325-333), 86 candidates, 0 transcription errors corrected,
%  86 the witness's fault, 0 unresolved.  See TRANSCRIPTION-PLAYBOOK.md
%  §3 step 4.  The transcribers took their words from that layer
%  corrected at the image, so the diff cannot see a layer error they
%  copied; the second reading exists for that class.
%  SECOND READING of all 41 pages, word by word and mark by mark at the
%  image (300 dpi; 600 dpi for any doubt), by five readers who had
%  transcribed nothing and were forbidden the text layer: no word,
%  punctuation, emphasis or paragraph discrepancy in the text as
%  transcribed.  The readers settled the file's open points: p. 314 the
%  mark over „eller“ is a printed acute, so the word now stands as
%  printed, „ellér“ (confirmed at 600 dpi in the caller); p. 318 the final
%  letter of „tidligere“ is a wrong sort (now logged as a printer's
%  error); p. 327 „theoride“ confirmed as printed.  Also noted: p. 310 the
%  final e of „hele“ is a raised sort (reading unaffected).  Result: no
%  transcription word error in 41 pages; one printer's error (p. 314) now
%  transcribed as printed instead of normalized.
% =====================================================================
\documentclass[12pt,a4paper,openany]{book}
\usepackage[utf8]{inputenc}
\usepackage[T1]{fontenc}
\usepackage{textalpha}
\usepackage{libertinus}
\usepackage{libertinust1math}
\usepackage{graphicx}    % for \reflectbox (the old Danish ɔ: id-est mark)
\DeclareUnicodeCharacter{0254}{\reflectbox{c}}  % ɔ (U+0254) = reversed c
\usepackage[danish]{babel}
\usepackage[protrusion=true,expansion=true]{microtype}
\usepackage{setspace}
\usepackage[margin=1.2in]{geometry}
\usepackage{fancyhdr}
\setlength{\headheight}{28pt}
\addtolength{\topmargin}{-13.5pt}

% Footnotes as the 1871 printing sets them: *).
\renewcommand{\thefootnote}{*)}

% The article's numbered sections: the number alone, centred on its line.
\newcommand{\secnum}[1]{\par\begin{center}#1\end{center}\par\nobreak}

\usepackage{hyperref}
\hypersetup{
  pdftitle={Positivisme og Materialisme},
  pdfauthor={Harald Høffding},
  hidelinks
}

\pagestyle{fancy}
\fancyhf{}
\fancyhead[LE,RO]{\thepage}
\fancyhead[RE]{\textit{Harald Høffding: Positivisme og Materialisme}}
\fancyhead[LO]{\textit{Nyt dansk Maanedsskrift, 2.\ Bind}}

\setstretch{1.3}

\title{\textbf{Positivisme og Materialisme}}
\author{Harald Høffding}
\date{Nyt dansk Maanedsskrift, udg.\ af Vilhelm Møller,
  2.\ Bind (April--September 1871), pp.~193--208, 309--333}

\usepackage{marginnote}
\newcommand{\opage}[1]{\marginnote{\footnotesize\textit{[#1]}}}
\newcommand{\apage}[1]{\marginnote{\footnotesize\textit{[#1]}}}

\begin{document}
\maketitle
\thispagestyle{empty}
\clearpage

% ===== Part 1 (Nyt dansk Maanedsskrift 2, pp. 193--208) =====
% --- p. 193 ---
% p. 193: the upper part of the page is blank (the article opens a new page; no running head, no folio -- the „13“ at the foot is the sheet signature, not transcribed); the head „Positivisme og Materialisme.“ in display blackletter (textura) on one line; below it the subtitle in three lines in a heavier (semi-bold) Fraktur of about body size, „Maudsleys“ letterspaced (gaps 10, 10, 11, 9, 8 px), broken „Nerve=|systemet“, footnote mark „*)“ after „1870)“; a short centred rule (c. 1.8 cm); the numeral „1.“ centred in roman type; the text opens with an enlarged Fraktur initial D, no indent; a short footnote rule at the left above the footnote
\opage{193}\begin{center}\textbf{\large Positivisme og Materialisme.}\end{center}

% footnote *) on p. 193
{\small Med særligt Hensyn til \emph{Maudsleys} Foredrag: „Om Forholdet mellem Sjæl og Legeme og mellem Sindssygdomme og andre Lidelser af Nervesystemet“. (Oversat i „Ugeskrift for Læger“. Septbr.-Oktbr. 1870)\footnote{Senere udkommet i særligt Aftryk.}.\par}

\begin{center}\rule{1.8cm}{0.4pt}\end{center}

\secnum{1.}

Det er først i den nyere Tid, man kan tale om Videnskaber.
Oldtiden og Middelalderen kjende egenlig kun een,
altomfattende Videnskab. Aarsagen hertil ligger i det væsenlig
forskjellige Synspunkt, hvorfra Tilværelsen kan opfattes.
Det ene Hovedsynspunkt er Heelhedens, det ene, sidste Formaal,
hvortil Alt samler sig for den begribende Tanke, og
som derfor maa være Slutnings- og Høidepunktet i enhver
Verdensanskuelse. Oldtiden og Middelalderen anlægge ene
% emphasis: „dette“ letterspaced (gaps 10, 8, 10, 9 px against 1--5 px solid)
\emph{dette} Synspunkt, og derfor bliver Videnskaben for dem kun
een. Af Heelhedens Idee, der i Oldtiden væsenlig sees som
det Skjønnes Idee, i Middelalderen som det Godes Idee,
søger man saa at udlede det hele System af Betingelser, som
danner Grundlaget for Ideens Virkeliggjørelse; Betingelserne
sees som indeholdte i Ideen selv, uden at der føles Trang til
en særegen Forsken med særegne Love og Methoder for at
undersøge dem. I den nyere Tid derimod har Tanken
særlig rettet sig mod Betingelserne som saadanne; den har
% --- p. 194 ---
\opage{194}indseet Umuligheden af at naae til deres Erkjendelse uden
gjennem en eiendommelig, specielt med dem syslende Forsken,
og lidt efter lidt løste sig derfor den ene Kreds af Betingelser
efter den anden fra den store Heelhed som Gjenstand for en
enkelt, afsluttet Videnskab med egen Methode og egne Principer.
Erfaringen indsættes i sin Ret og bringes under
exakte Love.

Denne Opløsning af Videnskabens abstrakte Eenhed
kunde ikke gaae for sig uden Modstræben fra deres Side,
som med denne absolute Eenhed troede at miste al videnskabelig
Eenhed i Erkjendelsen. Vel var det ved den nyere
Tids Begyndelse især de middelalderlige Aristotelikere, der modarbeidede
Erfaringsvidenskaben; og for dem var Eenhedsidealet
stivnet til en Autoritetstro, der lod dem sætte de
% footnote *) on p. 194
% emphasis in the footnote: „Galilæi“ spaced antiqua (gaps 8, 9, 9, 9 px); „Feuerbach“ spaced Fraktur (gaps 8, 9, 9 px); „(Dialogus de systemate mundi“ and the Latin quotation „Eqvidem ... fuisses!“ in antiqua; the dot under the s of „obstaret“ is a speck
aristoteliske Udsagn over deres egne Øines Vidnesbyrd\footnote{\emph{\textit{Galilæi}} \textit{(Dialogus de systemate mundi} anført af \emph{Feuerbach}: Pierre Bayle. 2te Ausg. S. 213) fortæller, at da en berømt Læge i Venedig paa en Forelæsning havde paaviist, at den største Nervestamme udgik fra Hjernen, spurgte han en tilstedeværende Aristoteliker, om han endnu troede, at Nerverne udsprang i Hjertet og ikke i Hjernen; denne svarede da, at han vilde have erklæret sig overbeviist ved, hvad han nu saae for sine Øine, hvis Aristoteles ikke udtrykkelig havde lært det Modsatte: \textit{Eqvidem ita aperte rem oculis subjecisti, ut, nisi textus Aristotelicus aperte nervos ex corde deducens obstaret, in sententiam tuam protracturus me fuisses!}}.
Men den spekulative Idealisme i Tydskland forsøgte i dette
Aarhundredes Begyndelse atter at udvikle Alt af et eneste
Udgangspunkt, Ideen. Schelling erklærede, at Baco havde
fordærvet Filosofien, Newton Fysiken! Intet har dog
kunnet standse Erfaringsvidenskabens seirrige Fremgang;
Skridt for Skridt er den trængt frem og har omspændt Alt
med den exakte Methodes Net.

Denne store Aandsmagt har ikke undgaaet den Fare,
der, efter hvad de Gamle saa dybsindigt lære os, truer enhver
lykkelig og seirende Magt: at overskride sin Grændse og sin
% --- p. 195 ---
\opage{195}Berettigelse. Men det maa fremhæves, at Grændsen her
ikke er let at drage. Om ydre, endelige Grændser kan der
paa dette Omraade ikke være Tale. Grændsen maa ligge
oprindelig i selve det Princip, hvorfra der gaaes ud. Den
moderne Erfaringsvidenskab spørger ikke om Gjenstandenes
ideelle Betydning og Værd; den søger kun at bestemme deres
Elementer og indre Sammenhæng; den gaaer ud paa et
% emphasis: „hvorledes“ (gaps 10, 9, 9, 11, 9, 9, 9, 10 px) and „hvorfor“ (gaps 13, 10, 8, 9, 8, 8 px) letterspaced
\emph{hvorledes} og ikke paa et \emph{hvorfor}. Men saa følger det
da af sig selv, at Resultatet ogsaa kun bliver en Erkjendelse
af Elementerne og deres indbyrdes Forbindelse, --- at der
altsaa til en total, afsluttende Erkjendelse af Tilværelsen udfordres
som Supplement en Erkjendelse, der fastholder det
ideelle Synspunkt og gaaer ud fra den totale Eenhed. Kun
Vexelvirkningen og Sammenslutningen af disse Synspunkter
kan føre til Maalet; det er de to Ender, der maae forbindes,
for at Kredsen kan sluttes. Erfaringsvidenskaben overskrider
altsaa sin Grændse, naar den umiddelbart mener at have
% emphasis: „hele“ letterspaced (gaps 10, 8, 8 px)
løst det \emph{hele} Tilværelsesproblem. Den holder kun den ene
af Traadene i sin Haand.

Paa dette Punkt ligger Spiren til Materialismen.
Ideen, den totale Eenhed er Formens Bestemmelse, overfor
hvilken Betingelsernes System repræsenterer Stoffets, Materiens
Bestemmelse, og Materialismens Eiendommelighed
bestaaer nu i at oversee Formens, Ideens Selvstændighed og
ene at udlede den af Stoffets Sammenhæng. Den betragter
altsaa Stoffet ikke blot som Betingelse for Formen, men
som Formens Aarsag. Den materialistiske Tendens ligger
saaledes for det ukritiske Blik Erfaringsvidenskaben saa nær,
at det intet Under er, at den i vore Dage tør ansee sig selv
som dennes eneste retmæssige Datter. Hos Naturforskere, der
ikke have et skarpt og sikkert Øie for deres Videnskabs
Væsen, indfinder Materialismen sig let, saasnart de fra
Specialundersøgelsernes Detail søge at hæve sig til en mere
omfattende Totalanskuelse (i Reglen i meget populær Form);
de glemme, at der til en saadan kommer andre Elementer
i Betragtning end de, den enkelte Videnskab kan omfatte.
% --- p. 196 ---
\opage{196}Dette hører til de uheldige Følger af Arbeidets Deling paa
Erkjendelsens Omraade.

Dog er det ikke blot, naar det gjælder om den sidste,
afsluttende Totalanskuelse, at den materialistiske Tendens,
saaledes som vi have karakteriseret den, kan gjøre sig gjældende.
Ogsaa indenfor Erfaringsvidenskaben selv gjentager
det samme Forhold sig mellem Betingelse og Aarsag, Stof
% emphasis: „der“ letterspaced (gaps 11, 9 px)
og Form; ogsaa \emph{der} er altsaa en Mulighed til modsatte
Retninger. Jo mere den nyere Videnskabs store Fremskridt
paa ethvert Omraade skyldes den exakte Opfattelse af Betingelserne,
des nærmere ligger den Misforstaaelse, at Sagen
selv ikke er Andet end Betingelserne. Materialismen er
Reaktionen mod tidligere Tiders Dualisme, en Protest mod
den udvortes Adskillelse af Væsenet fra dets Betingelser.
Hvad den bekæmper, er egenlig blot den Mening, at der
ligesom bag ved Fænomenerne skulde være skjult en vis
mystisk Kraft, som var uafhængig af Betingelserne. Den
stræber derfor at føre Alt tilbage til det Elementære og lade
det opstaae af dettes Sammensætning. Den høiere, mere
konkrete Videnskab maa da stedse kæmpe mod Overgreb
fra den lavere, mere abstrakte og elementære. I Fysiken
betegnes den Retning som materialistisk, der nægter de enkelte
mekaniske Kræfters Eiendommelighed og kun betragter dem
som specielle Former af elementær Bevægelse. Udstrækkes
dette videre, vil ikke blot Kemien opløses til mekanisk Fysik,
men ogsaa Fysiologien til Kemi, Psykologien til Fysiologi.
Det Konkrete og Eiendommelige skal paa ethvert Punkt kun
sees som Produkt af de abstrakte, almene Betingelser, den
individuelle Eenhed kun som Sammensætning af mekanisk
virkende Elementer.

% emphasis: „Tendens“ letterspaced (gaps 11, 10, 11 px)
Vi tale her alene om den materialistiske \emph{Tendens}.
Fuldstændig gjennemføres kan Materialismen ikke. Den rene
Materie vilde være Et med det rene Intet, hvoraf Intet kan
komme. Materialismens Væsen ligger derimod i det Relative,
som jo Stoffets Bestemmelse bestandig er relativ, idet
det, som i den ene Henseende er Stof, i den anden er Form.
% --- p. 197 ---
\opage{197}Den relative Materialismes Berettigelse ligger ifølge det
Foregaaende i at udtrykke den Sandhed, at Stoffet forklarer
% emphasis: „forsaavidt“ letterspaced (gaps 13, 11, 10, 11, 10, 11, 12, 10, 12 px)
Formen, \emph{forsaavidt} denne beroer paa og afhænger af hiin.
Men naar Tendensen gaaer videre, træder Selvmodsigelsen
klart for Dagen.

Thi Formen, Heelhedens Idee kommer ikke til udenfra,
men lægger sig for Dagen indenfor selve Betingelsernes
Verden. Den finder sit særlige Udtryk --- skjøndt den er
overalt tilstede --- hvor en eiendommelig Kreds af Love og
Fænomener udfolder sig paa Grundlag af mere abstrakte og
elementære Forudsætninger, og hvor derfor en ny og høiere
Videnskabs Omraade begynder. Disse forskjellige Kredse
danne ligesom en Knudelinie, hvor de kvantitative Overgange
fortætte sig til kvalitative Midtpunkter, fra hvilke Bevægelsen
gjennem stadig Kvantiteren stedse gaaer videre til
% footnote *) on p. 197, continued at the foot of p. 198 (from „punkter ikke absolut“); given whole here at its mark
% emphasis in the footnote: „Lewes“ letterspaced (gaps 10, 8, 8, 10 px)
andre og atter andre Knudepunkter\footnote{Smlgn. hermed hvad \emph{Lewes} i sin Biografi af Goethe siger om dennes Farvelære. Den er forfeilet som et Bidrag til Lyslæren, fordi den seer bort fra den exakte kvantitative Methode. „Uden Vægtskaalens frie Kontrol kan det kemiske Experiment aldrig blive kvantitativt, og uden kvantitativ Kundskab kan der ikke være nogen kemisk Videnskab strengt taget, men kun kvalitativ ɔ: tilnærmelsesviis Kundskab. Ingen Grad af Iagttagelse kan gjøre Iagttagelsen nøiagtig, naar den ikke kan maales. Om man saa iagttager faldende Legemer i al Evighed, uden Mathematik kommer man aldrig til Gravitationsloven.“ Men paa den anden Side er Goethes Farvelære meget værdifuld, naar man kun seer hen til Opfattelsen af Farverne som saadanne, uden Hensyn til deres mekaniske Betingelser. „For Kunstnere og Fysiologer, det vil sige for dem, som navnlig have Interesse for Farvefænomenerne som Sandseindtryk, og som mere fordre kvalitativ end kvantitativ Kundskab, have hans Arbeider stort Værd“. --- Hertil er at bemærke, at det kun er fra et abstrakt kvantitativt Synspunkt, at det Kvalitative kan gjøres til Et med det Tilnærmelsesvise. Efter Lewes's egen Karakteristik betegner det Kvalitative derimod Fænomenet som Heelhed og den Betydning, det Værd, det har som saadan, og dette bliver her, hvor der er Tale om Farverne, Et med den ideelle, psykologiske og æsthetiske Betydning. Naturligviis kunne de forskjellige Synspunkter ikke absolut holdes ude fra hinanden; det er netop dette, der betegnes ved Udtrykket Knudelinie.}. Materialismen
% --- p. 198 ---
\opage{198}grunder sig nu her paa et optisk Bedrag, idet den paa
Grund af, at det Kvalitative kan bestemmes gjennem Kvantiteren,
mener, at det i sig selv kun beroer paa Størrelsesforskjelle.
Men da hine Knudepunkter faktisk existere i
Tilværelsen, kommer den materialistiske Tendens stedse til at
modsige sit Princip, kommer til at indsmugle kvalitative
Bestemmelser, medens den mener at forklare Alt kvantitativt.

Theoretisk seet er Materialismen i vor Tid altsaa let
at forklare; den hænger sammen med, hvad der fra en anden
Side seet udgjør vor Tids Storhed og Styrke. Den er et
Tidsfænomen, hvis Væsen og Oprindelse man kun behøver
at gjøre sig klar for at kunne overvinde det. I Praxis
fremtræder Materialismen vel med større Seighed og Styrke;
men dog maa ogsaa her det Samme gjælde. Fra theologisk
Side er der udtalt en umiddelbar Forkjættrelse over Materialismen,
idet denne ikke betragtes som en Vildfarelse i Erkjendelsen,
men som en af den onde Villie til sin egen
% footnote *) on p. 198, continued at the foot of p. 199 (from „die Abtreibung“); given whole here at its mark
% the German quotations in the footnote are in Fraktur, not antiqua; only „Dr.“ is antiqua
% as printed in the quoted German (p. 198): „mekanischen“ (k) and „Rechtsbewußtseins“ (p. 199, with -s); not corrected
Besmykkelse uddannet Lære\footnote{Smlgn. „Briefe gegen den Materialismus von \textit{Dr.} F. Fabri, Missionsinspektor“. 2te Aufl. S. 216: „Die Lehre der Materialisten von einer mekanischen Denknothwendigkeit ist im Grunde nichts als das Feigenblatt ihrer die moralische Verantwortlichkeit läugnenden Unfreiheit des Willens. Der verkehrte Wille vollendet sich in der verkehrten Theorie“. Forfatteren mener nemlig, at al Viden gaaer ud fra en umiddelbar Tro, som igjen bestemmes ved Menneskets Villie. Derved gjør han egenlig sin hele Bog overflødig, medens det tillige bliver konsekvent, naar han tillægger sine Modstandere uethiske Motiver, hvilket han driver til fanatisk Yderlighed. Han anfører saaledes i en Note, at en „Verein badischer Aerzte zur Förderung der Staatsarzneikunde“ havde udnævnt Büchner, Forfatter til det bekjendte materialistiske Skrift „Kraft und Stoff“, til sit Æresmedlem. Det vilde, tilføier han (S. 27), ikke være uinteressant at faae Noget at vide om de Motiver, der have ledet Foreningen for „Staatsarzneikunde“ ved hiint Skridt: „Doch wir erinnern uns so eben, daß der Verfasser von „Kraft und Stoff“ die Abtreibung der Leibesfrucht als ein Recht der Mutter wieder vindicirt hat, und daß er die Freiheit des Willens und also die sittliche Zurechnungsfähigkeit mit Herrn Vogt läugnet! Es kann aber nicht in Abrede gestellt werden, daß dies Behauptungen sind von solcher Tragweite, die allerdings nicht nur die gesammte gerichtliche Medicin, sondern unser ganzes Rechtsbewußtseins und alle sociale Ordnungen in Grund und Boden reformiren müßten. Und bei der Lage des Falls bleibt wenigstens bis zur Führung des Gegenbeweises (!) nichts übrig, als anzunehmen, jener Carlsruher Verein habe in Herrn Büchner den verkannten Apostel einer solchen schönen Zukunft ehren wollen“. Materialisterne toge ikke læmpeligt paa deres Modstandere, men man vilde i deres Skrifter forgjæves søge saadanne giftige Insinuationer som den her anførte. Jeg behøver neppe at bemærke, at Büchner aldeles ikke forsvarer Fosterfordrivelse; han anfører kun, for at vise de moralske Forestillingers Vexlen til forskjellige Tider, at medens vi ansee den for Brøde, fandt Romerne intet Umoralsk deri. Det er Fabri selv, der i sin Iver for den „Mission“, han har paataget sig mod Materialismen, heri finder en „Hævdelse af Fosterfordrivelse som en Ret for Moderen“!}. For den filosofiske Kritik er
% --- p. 199 ---
\opage{199}det praktiske Korrektiv derimod i nøie Sammenhæng med det
theoretiske og fremtræder ikke isoleret. Derved bliver den
ogsaa istand til at anerkjende det i den materialistiske Tendens
Berettigede.

Der er i vore Dage gjort et storartet Forsøg paa at
lægge det Berettigede i den materialistiske Tendens til Grund
for en heel Verdensanskuelse. Den saakaldte Positivisme,
% emphasis: „Auguste Comte“ letterspaced (gaps 9, 11, 8 and 8, 8, 9 px)
hvis Stifter er \emph{Auguste Comte}, og som har vundet en
stor Tilslutning blandt de franske og engelske Filosofer og
Naturforskere, lægger an paa at gjennemføre en paa Erfaringsvidenskabens
Princip umiddelbart grundet Filosofi, ---
en „positiv“ Filosofi, fordi den er støttet paa de positive
Videnskaber. En Betragtning af denne Retning vil klarere
end almindelige Overveielser godtgjøre den relative Materialismes
% emphasis: „Maudsleys“ letterspaced (gaps 10, 11, 10, 11, 14, 9, 9, 10 px)
Eiendommelighed. \emph{Maudsleys} interessante Foredrag,
der med Rette have vakt Opmærksomhed i en vid
Kreds, høre, hvad den hele til Grund liggende Tankegang
angaaer, ind under den almindelige positivistiske Retning.
% p. 199: section 1 ends at the foot of the text column with no rule; the faint short rule below the footnote is show-through/set-off (it sits under a ghost of the line „omvendt at forklare Mennesket ved Verden.“ and prints at about half the density of the footnote rule), not transcribed
% --- p. 200 ---
% p. 200: the numeral „2.“ centred in roman type directly under the running head; no rule before it
\opage{200}\secnum{2.}

Auguste Comte stiller sin positive Filosofi i Modsætning
til tidligere Tiders theologiske og metafysiske Filosoferen.
Paa det første, ufuldkomne Bevidsthedstrin seer Mennesket
Alt i Naturen fremgaaet af en lignende Aarsag som den, der
bevirker hans egne Bevægelser og Handlinger, altsaa af en
bag ved Fænomenerne liggende Tanke og Villie. Da han
endnu ikke har nogen klar Bevidsthed om nødvendig Lovbestemthed,
seer han den i Naturen virkende Villie som en vilkaarlig
Villie, der skalter og valter efter Behag med Tingene; og da
dette igjen forudsætter en høiere Magt end den menneskelige,
som ikke har absolut Herredømme over Naturen, saa bliver
hiin vilkaarlige Villie for Mennesket den guddommelige Almagt,
til hvilken han forholder sig gjennem Religionen.
Men ved den nøiere, vedvarende Beskjæftigelse med Fænomenerne
opstaaer efterhaanden den Erkjendelse, at der ikke
hersker absolut Vilkaarlighed, men fast Orden og Regelmæssighed
i deres Løb. Istedetfor den almægtige, personlige
Villie fremtræder saa for Bevidstheden Tingenes indre Væsen
som den egenlige Aarsag. Bag Fænomenerne tænker man
sig skjulte Kræfter, i hvilke den sidste Forklaring af dem
søges; enhver af disse Kræfter raader i sin Kreds af Fænomener,
og Systemet af alle Kræfterne danner en indre
Verden, som frembringer og opholder den ydre, synlige
Verden. Istedetfor hiin Verden af Guder, i hvilken den
umiddelbare Bevidsthed søgte Fænomenernes Forklaring, er
da her sat en Verden af metafysiske Væsener \textit{(entités métaphysiques)}.
Det er ikke længer Dryader, som besjæle
Træerne og lede deres Væxt og Liv, men Aarsagen til
Plante- saavelsom til Dyrelivet søges i en „vegetativ Sjæl“
eller en „Livskraft“. Analogien med den menneskelige Bevidsthed
træder her mere tilbage. Men dog er det fælles
for det religiøse og det metafysiske Bevidsthedstrin, at de
begge forklare det Ydre ved det Indre, at de lægge Menneskets
Væsen til Grund for Forklaringen af Verden, istedetfor
omvendt at forklare Mennesket ved Verden.

% --- p. 201 ---
% p. 201: p. 200 ends its paragraph („ved Verden.“) at the foot with no footnote, so nothing carries over; p. 201 opens with an indented new paragraph, the Comte quotation set in antiqua with roman guillemets »...« (kept as printed, outside \textit, as in the house form for „pagani“)
% the French is printed thus: „caracterisé“ (for caractérisée), „rationelle“, „2éme“ -- transcribed as printed
\opage{201}»\textit{Le véritable esprit général de toute philosophie théologique
ou métaphysique consiste à prendre pour principe,
dans l'explication des phénomènes du monde extérieur,
notre sentiment immédiat des phénomènes humains; tandis
que, au contraire, la philosophie positive est toujours caracterisé,
non moins profondément, par la subordination
nécessaire et rationelle de la conception de l'homme à celle
du monde}« \textit{(Auguste Comte: Cours de philosophie positive.
2éme éd. Tome III. p. 188)}.

Den tredie Bevidsthedsperiode betegnes ved den positive,
videnskabelige Methodes Herredømme. Kun fordi den religiøse
og den metafysiske Bevidsthed gik ud fra Menneskets
Væsen, kunde de mene at naae en Erkjendelse af Tingenes
Væsen og Aarsager, en Indsigt i Verdens Oprindelse og
Bestemmelse. Herpaa giver Tanken Afkald, naar den er
naaet til »\textit{l'état positif}«. Den holder sig nu til det Positive,
% printer's error: the quotation „positif, ... constatés“ is closed with » (an opening guillemet turned the wrong way) instead of «; transcribed as printed
det vil sige det faktisk Givne: »\textit{positif, c'est à dire
fondé sur des faits bien constatés}». Men det faktisk Givne
er kun Fænomenerne. Tankens eneste Opgave er da at
finde deres virkelige Love, det vil sige, det uforanderlige Forhold,
i hvilket de følge efter hverandre og ligne hverandre.
Derimod er det umuligt at kjende deres Natur og Væsen i
og for sig. Væsenet i og for sig vilde være det Ubetingede
og Absolute; men det Positive er stedse betinget og relativt.
Den positive Filosofi maa derfor opstille som sin absolute
Grundsætning, at vi intet Absolut kunne erkjende.

Hvad Comte kalder »\textit{la loi des trois états}« fører os altsaa
ind i Positivismens eiendommelige Væsen. Men nærmere Bestemmelser
vise sig at være nødvendige. Vi erkjende kun det Positive,
ɔ: (saaledes som Comte bestemmer det) det Relative, det Betingede.
Men vil dette sige, at kun det Relative er, at der intet
% emphasis: „gives“ letterspaced (gaps 13, 11, 14, 10 px); „erkjende“ letterspaced (gaps 11, 16, 12, 13, 10, 11, 12 px)
absolut Princip \emph{gives}, eller vil det kun sige, at vi ikke kunne
\emph{erkjende} det? Naar Alt, hvad vi erkjende, er betinget, skal
saa det sige, at Alt kun er Betingelser, saa at det Betingede
kun er Betingelsernes afledede Egenskab? Der er med andre
% --- p. 202 ---
\opage{202}Ord i Positivismens Princip en Mulighed saavel til Skepticisme
som til Materialisme.

Da Comte afviser al Erkjendelse af de indre Aarsager,
af Tingenes Væsen og Princip, gaaer hans Stræben kun
ud paa at opdage Fænomenernes Love, ɔ: de Betingelser,
under hvilke de opstaae. Disse Love eller „almindelige
Kjendsgjerninger“ \textit{(faits généraux)}, som gjælde for en særegen
Kreds af Fænomener, vise saa igjen tilbage til mere
almindelige Forhold, af hvilke de kun ere en særegen Form,
og ved saaledes stedse at gaae tilbage kommer man tilsidst
fra de konkrete, indviklede og sammensatte Fænomener til de
simpleste, meest elementære og derfor meest almindelige Love.
De mathematiske Love omfatte alle Sfærer i Tilværelsen,
fordi de ene forudsætte de rene Størrelsesforhold. »\textit{La loi
des trois états}« findes her fuldstændig realiseret: Mathematiken
har allerede længe været en positiv Videnskab. Derimod
gjælder dette ikke om de andre Videnskaber. Her forudsættes
endnu stedse særegne Principer, som ikke reduceres
til simple positive Fænomener eller Forhold mellem Fænomener.
% emphasis: „da“ letterspaced (d--a gap 10 px against 4--5 px in „paa“ beside it)
Positivismens Fuldendelse beroer \emph{da} paa, at dette
skeer, at enhver Eiendommelighed, ethvert særeget Princip
føres tilbage til elementære Forudsætninger, og at disse atter
sees fremgaae af den forudgaaende, mere abstrakte og derfor
mere omfattende Kreds af Love. Vel forudsætter enhver
saadan Kreds stedse en egen Erfaringsmethode; men den
kan ikke blive systematisk uden ved at udledes fra de mere
elementære Kredse.

Hvad der her antydes, gaaer i Virkeligheden ud over,
hvad der indeholdes i Positivismens oprindelige Program.
Positivismen maa nødvendigviis ophøre, hvor det „Positive“
hører op, og dette skeer, saasnart den faktiske Forbindelse af
Fænomenerne hører op. Den Lov, Comte mener at have
fundet, at de almindeligste Egenskaber, som gjælde i det
videste Omfang, tillige ere de meest elementære og abstrakte, ---
en Lov, der, som man med Rette har bemærket, kun er den
velbekjendte logiske Sætning om det omvendte Forhold mellem
% --- p. 203 ---
% footnote *) on p. 203
\opage{203}et Begrebs Indhold og dets Omfang\footnote{Et Begrebs Indhold er de Kjendemærker, det angiver, dets Omfang er Alt, hvorpaa disse Kjendemærker passe, og som derfor hører ind under Begrebet. Jo større nu Indholdet er (jo flere Kjendemærker, Begrebet angiver), des mindre er Omfanget (de Individer, som besidde Begrebets Kjendemærker og altsaa kunne henføres under det). F. Ex. Begrebet Dyr --- Begrebet Hest: dette indeholder flere Kjendemærker, men omfatter færre Individer end hiint.} --- denne Lov indeholder
ingen Borgen for, at den reale („positive“) Forklaring
af Fænomenerne skulde kunne findes ved stedse at gaae tilbage
fra det Konkrete og Kombinerede til det Abstrakte og
Elementære. Comte har her uden Videre overført et logisk
Forhold paa det Reale og Positive og har derved stillet sig
paa et lignende Standpunkt som den spekulative Filosofi, der
udleder det Reale af det Logiske. Han forklarer her, for at
bruge hans egne Ord, Verden ved Mennesket, ikke Mennesket
ved Verden.

Den Opgave, han har stillet sig, er at udvikle de
Grundbegreber, paa hvilke al vor Tænkning og Erkjendelse
beroer, og han erklærer udtrykkelig, at disse Grundbegreber
ikke selv ere grundede paa noget „forudgaaende
intellektuelt System“. Men i den Anvendelse, han gjør af
hiin logiske Lov, har han en ubevidst Forudsætning, som ikke
inddrages i Undersøgelsen. Her viser Positivismen sig da
som Dogmatisme og gaaer herved tillige over i Materialismen.
Har man indrømmet, at den logiske Sætning om
det omvendte Forhold mellem Begrebets Indhold og Omfang
umiddelbart kan anvendes ved Forklaringen af Naturen,
saa har man derved indrømmet selve det materialistiske Princip.
Men ved denne umiddelbare Anvendelse oversees, at
Sætningen undergaaer en betydelig Forandring ved at overføres
paa Naturens Omraade paa denne Maade. Logisk seet
er det abstrakte Begreb det høiere, forsaavidt det indbefatter
de konkretere Begreber i sig; men naar man ikke bliver
staaende ved en reent formel Betragtningsmaade, maa det
% --- p. 204 ---
\opage{204}høiere Begreb udvikles saaledes, at det ikke blot fremkommer
ved Subtraktion af de enkelte Karakteermærker, men indeholder
disse som Muligheder, som Momenter i sig, saaledes at det
giver den sidste, afsluttende Forklaring af det Konkrete, overfor
hvilket det har det Fortrin ikke at være bundet til det
Enkelte og Endelige, og dog ideelt seet at indbefatte dette i
sig. Tilbagegangen fra det høiere, nu ikke abstraktere men
konkretere Begreb til det lavere, mere begrændsede skeer da
ikke ved en blot Addition af de enkelte specifike Karakteermærker,
men disse udvikles af hiint som særlige, relative
Udtryk derfor. Saaledes som Comte opfatter det høiere Begreb
og anvender det ved Fænomenernes Forklaring, er det
det blot abstrakte, ved Fradrag af det Særlige udtømte
Almeenbegreb. Han betragter saaledes det Organiske blot
som et særligt Tilfælde af det Kemiske o. s. v. --- ethvert
Fænomen i sidste Instants som et særligt Tilfælde af det
Mathematiske. Men hvorledes kan han da gjøre Overgangen
fra det Almindelige til det Særlige? Strax ved Overgangen
fra Mathematiken til Fysiken kommer Spørgsmaalet om
Materien frem, et Spørgsmaal, som heelt omgaaes. Paa
samme Maade forudsættes bestandig det Særlige, medens det
efter hans Lære stedse skulde vises som et af de abstrakte
% printer's error (or oddity): „Sammensat“ printed with a capital S in mid-sentence; transcribed as printed
Elementer Sammensat. Fordi Mathematikens Love ere gyldige
i hele Tilværelsen, ere alle Fænomenerne derfor ikke
blotte mathematiske Exempler; fordi Kemiens Love ikke ophæves,
men tvertimod bekræftes ved Undersøgelsen af de
organiske Væsener, blive disse ikke til blotte kemiske Produkter.
Den fundamentale Feil er her stedse Forvexlingen af Aarsag
og elementær Betingelse.

Ved sit oprindelige Princip er Positivismen væsenlig
Empirisme. Det Positive bestemtes jo saaledes, at det faldt
sammen med det faktisk Givne, det Fænomenale, det Relative.
Auguste Comte erklærer derfor Materialismen for
uvidenskabelig, fordi den opstiller et Absolut som sidste Grund
til Alt og derved gaaer ud over det Givne. Ved hele den
Modsætning, i hvilken han stiller Positivismen som den, der blot
% --- p. 205 ---
\opage{205}gaaer ud paa at konstatere Fænomenernes empiriske Sammenhæng,
til Metafysiken som den, der søger de indre Grunde,
Tingenes Væsen i og for sig, nærmer Comte sig den skeptiske
Empirisme, saaledes som Hume har grundlagt den. Det
var netop fra Nægtelsen af Aarsagsbegrebets Gyldighed,
Hume gik ud, og Erkjendelsen blev for ham en saadan
Forbindelse af Fakta, som Comte betegner ved positiv Filosofi.
Den bekjendte Fornyer af Humes Lære i vor Tid, John
Stuart Mill, var en Discipel af Comte, hvis Anskuelse han
opfattede og udviklede i empirisk-skeptisk Retning. Men Mill
regner det som en af Comtes Feil, at han aldrig indrømmer
aabne Spørgsmaal. Navnlig mener han, at fordi den
positive Viden kun kan vise, at i Universet, eller rettere den
Deel af Universet, som vi kjende, ere Fænomenerne knyttede
til hinanden paa en bestemt Maade, derfor kan den ikke nægte
Muligheden af, at hele denne lovbestemte Tilværelse har sin
sidste Grund i en evig Tanke, naar man kun ikke lader
denne virke paa vilkaarlig og lovløs Maade. Hvad Mening
man danner sig herom, beroer ifølge Mill paa den Betydning,
man tilskriver deels de Analogier, der kunne findes i
Naturen med den bevidste Stræben efter Formaal, deels
% footnote *) on p. 205; the footnote is set wholly in antiqua
Menneskeslægtens almindelige Traditioner\footnote{\textit{Auguste Comte and Positivism, p. 15.}}. Der er her
indført en Skepsis, der vender sig mod selve Positivismens
Princip og nedsætter dette til kun at have relativ og betinget
Gyldighed. Men naar Positivismen ikke længer kan fastholde
sin egen absolute Gyldighed, saa er ogsaa den Skranke
falden, der adskilte den fra Skepticismen.

Mill finder Forklaringen af visse extravagante Lærdomme
i Comtes senere Skrifter i en Tendens til Eenhed og
Systematiseren, som efter hans Mening er eiendommelig for
franske Tænkere. Denne Tendens gjør sig gjældende i
Comtes hele Grundanskuelse; det var den, der førte ham til
dogmatisk at opstille Sætningen om det Konkretes Forklaring
% --- p. 206 ---
\opage{206}ved det Elementære, hvorved han gik over i Materialismen.
Denne Side af hans Lære blev meest fremhævet af hans
franske Tilhængere, navnlig af den bekjendteste blandt dem,
den lærde Natur- og Sprogforsker Littré. Stuart Mill
skiller sig derimod fra ham ogsaa i den Maade, hvorpaa
den „positive“ Videns Maal og Væsen nærmere opfattes.
Medens Comte søger at ophæve alle særlige Fænomener ved
at udlede dem af deres Elementer, indskærper Mill, at de
sidste Naturlove, til hvilke man kan naae, ikke kunne være
mindre talrige end de kvalitativt forskjellige Fornemmelser og
% printer's error: the antiqua parenthesis carries a closing guillemet » (turned the wrong way) with no opening guillemet anywhere before it; transcribed as printed
Bevidsthedstilstande \textit{(sensations or other states of consciousness}»);
Grændsen maa sættes, hvor det Kvantitative og
% footnote *) on p. 206; the footnote is set wholly in antiqua
Gradvise hører op\footnote{\textit{System of logic. 5 ed. II. p. 5.}}. Ogsaa ved denne Bestemmelse af den
positive Videns Grændse vises vi i en ganske anden Retning
end den materialistiske.

Men heller ikke Auguste Comte selv fastholdt hiin materialistiske
Grundsætning. Da han i »\textit{Cours de philosophie
positive}« kom til at omhandle de organiske Væsener, opstillede
han en Bestemmelse af Organismens Væsen, der
minder om den kantiske. Ved de uorganiske Væsener, siger
han, gaaer Erkjendelsens Vei fra Delene til det Hele;
men ved de organiske Væsener maa den sande Erkjendelse
af Delene udledes af Erkjendelsen af det Hele.
% p. 206: a stray raised mark (spacing material or a broken sort) stands between „nemlig“ and its comma; not transcribed
Det Eiendommelige ved det organiske Væsen er nemlig, at
det er en sluttet Eenhed, der udvikler sig; Forklaringen af
det maa da vise Grunden til denne Eenhed og Udvikling.
Analysen, som opløser Heelheden i dens elementære Dele, er
her ikke tilstrækkelig; den analytiske Virksomhed er kun Forberedelsen
til den afsluttende Erkjendelse, som beroer paa en
Heelheden opfattende Synthese. Dette leder nu Comte til
at opstille en Sætning, der er hiin tidligere fuldstændig
modsat. Det Høiere, siger han nu, forklarer det Lavere, ikke
omvendt. At udlede og forklare det Høiere af det Lavere er
% --- p. 207 ---
\opage{207}Materialisme. Denne afviser han altsaa nu som uvidenskabelig,
ikke længer fordi den gaaer ud over det Empiriske,
men fordi den ene gaaer ud fra det Elementære.

Den nye Lov, Comte her opstiller, anvender han dog
ikke paa et Punkt, hvor dens Anvendelse netop synes at
ligge nær, paa Forholdet mellem Sjæl og Legeme. Stuart
Mill bebreider ham med Rette, at han ikke har nogen Plads
for en psykologisk Videnskab, men gaaer lige fra Fysiologien
over til Sociologien (Læren om det menneskelige Samfund).
Grunden hertil er, at han ikke anerkjender
et eiendommeligt sjæleligt Princip, men betragter det
Sjælelige blot som et Produkt af det Legemlige. Her forbliver
han sin tidligere Grundsætning tro: at udlede det
Høiere af det Lavere. De aandelige Fænomener ere paa
eengang mere sammensatte og mere specielle end de fysiologiske
og fysiske. Naar der trods denne indviklede Sammensætning
viser sig en sjælelig Eenhed, et Jeg, saa er dette kun et abstrakt
Begreb, der danner sig efterhaanden ved Følelsen af
den Harmoni og Ligevægt, der i den normale Tilstand danner
sig mellem de mangfoldige, af hverandre uafhængige Kræfter
% p. 207: all three French quotations on this page close with » (as the one on p. 201 does), not «; transcribed as printed
og Funktioner i den menneskelige Natur. »\textit{La notion du
moi, c'est à dire, du consensus universel de l'ensemble de
l'organisme}». Det er derfor forgjæves at forsøge at gaae
ud fra Jeg'et, da dette kun udtrykker »\textit{un état purement
fictif}»; det Sande, som ligger til Grund, er den elementære
Mangfoldighed.

Ogsaa her gjælder det dog at fastholde Eenhedens
Realitet ligesaavel som Betingelsernes. Hvad Comte betegner
som »\textit{purement fictif}» er selve det menneskelige Aandslivs
Grundeiendommelighed; den lader sig ikke afvise med et
Magtsprog. Det er ikke blot i psykologisk Henseende, at
Jegets Begreb, Bevidsthedens Eenhed er af afgjørende Betydning.
Den menneskelige Erkjendelses Love vise tilbage til
Bevidsthedens indre Organisation. Naar Positivismen betragter
Erkjendelsen som en blot Samlen og Sammenfatten
af faktiske Fænomener og deres indbyrdes Forhold, da er
% --- p. 208 ---
\opage{208}dette en reen mekanisk og udadgaaende Virksomhed, der med
Nødvendighed viser tilbage til en indre Grund, hvoraf det
Mekaniske betinges, og hvorfra Virksomheden gaaer ud.
Og denne indre Grund kan ikke være Andet end selve den
samlende og sammenfattende Tanke, hvis Love lægge sig for
Dagen under denne samlende og sammenfattende Virken.
At Tankens Lov i sig er Et med Virkelighedens Lov, er
den sidste Forudsætning, hvorfra al videnskabelig Forsken
maa gaae ud. Først da vinder Erkjendelsen indre Nødvendighed
og sand Eenhed. Positivismens --- vi kunne her
tilføie: og Materialismens --- Uvidenskabelighed bestaaer da i
sidste Instans i, at den ikke undersøger selve Videnskabens
første Principer, ikke opstiller en Erkjendelseslære. Men den
% emphasis: „kan“ letterspaced (gaps c. 9--12 px against 1--5 px in „heller“ beside it)
\emph{kan} heller ikke indlade sig paa disse Undersøgelser uden at
opgive sig selv; dens Væsen er umiddelbar Dogmatisme.

Den filosofiske Kritik har at værne om Videnskabens
universelle Love og Principer, navnlig mod Forsøg paa
umiddelbart at opstille som videnskabelig Totalanskuelse, hvad
der kun har Gyldighed i en enkelt Sfære, fra en enkelt
Videnskabs Synspunkt. Filosofien selv har til sine Tider
troet i Kraft af sin universelle Karakteer at kunne beherske
de særskilte Videnskaber; i vore Dage er den modsatte Tendens
herskende. Den filosofiske Kritik vender sig mod begge
Eensidigheder. Den gaaer, som Sokrates fordum, omkring
og spørger dem, „der, fordi de paa en skjøn Maade kunne
udøve deres egen Kunst, ogsaa mene at være visest i andre
Henseender“, hvorledes de da kunne gjøre Fyldest med Hensyn
til Videns og Videnskabens sidste Grunde og afsluttende
Resultater. For denne Domstol godtgjøre Positivismen og
Materialismen deres Uvederhæftighed, --- og det er for denne
Domstol vi ogsaa indstævne Forfatteren af de interessante
% p. 208: the dot between „Sjæl“ and „og“ is a speck, not a hyphen
Foredrag „om Forholdet mellem Sjæl og Legeme“; ogsaa
han „øver sin egen Kunst saa skjønt“, at han har meent at
kunne gjøre Fyldest ogsaa udover dens Grændser.

% p. 208: part 1 closes with „(Fortsættes.)“ in body Fraktur, indented at the left; on the next line, set right, the signature „Harald Høffding.“ in bold Fraktur; below it a short centred double rule (checked at the image in the caller); nothing of the next article on this page
\noindent\hspace*{1.5em}(Fortsættes.)\par
\begin{flushright}\textbf{Harald Høffding.}\end{flushright}
\begin{center}\rule{2.2cm}{0.4pt}\\[-10pt]\rule{2.2cm}{0.4pt}\end{center}

\clearpage
% ===== Part 2 (Nyt dansk Maanedsskrift 2, pp. 309--333) =====
% --- p. 309 ---
% p. 309: folio „309“ at the head right; no running head; the upper part of the page is blank (no tail of a preceding article on this leaf -- only faint show-through of p. 310, including its running head and text); the head „Positivisme og Materialisme.“ in display blackletter (textura), as on p. 193; below it a short centred rule (c. 0.85 cm, x 643--742); the English epigraph in five lines of small upright antiqua (slab-serif), set to the right, broken „corre-|sponds“, with „John Stuart Mill.“ set right below; a second short centred rule of the same length and position (x 643--745); the numeral „3.“ centred in roman type; the text opens with a slightly enlarged Fraktur initial D, no indent
\opage{309}\begin{center}\textbf{\large Positivisme og Materialisme.}\end{center}

\begin{center}\rule{1.8cm}{0.4pt}\end{center}

\begin{flushright}\begin{minipage}{0.6\textwidth}\small\textit{Imperfect as is the science of mind, I do not scruple to affirm, that it is in a considerably more advanced state than the portion of physiology which corresponds to it.}\\ \hfill\textit{John Stuart Mill.}\end{minipage}\end{flushright}

\begin{center}\rule{1.8cm}{0.4pt}\end{center}

\secnum{3.}

Det er ikke tilfældigt, at de fleste af Materialismens Ordførere
i vor Tid ere Læger og Fysiologer. Studiet af den
menneskelige Organisme fører til at see, hvor afgjørende en
Betydning de legemlige Organer have for det hele aandelige
Liv. La Mettrie blev Materialist, da han under en heftig
Feber havde havt Leilighed til paa sig selv at iagttage Blodstrømningernes
Virkning paa Hjernen; fra dette Øieblik betragtede
han Tænkningen som en blot Følge af „den legemlige Maskines“
Organisation. Ifølge \textit{Système de la nature} løser
Medicinen det menneskelige Hjertes Gaade. Man har derfor
kaldet Medicinen Materialismens Theologi.

Dog er det langt fra, at det fysiologiske Studium
med Nødvendighed skulde føre til materialistiske Konsekvenser.
Man behøver kun at see hen til Mænd som Flourens og
Claude Bernard, Johannes Müller og Virchow for at forstaae
dette. Men i Almindelighed vil man fra fysiologisk
% --- p. 310 ---
\opage{310}Side udtale sig i positivistisk Retning. At særlig Maudsley
staaer paa Positivismens Standpunkt, vil let kunne vises.
Vel paakalder han kun et enkelt Sted udtrykkelig „Positivismens
Aand“; men at hans hele Betragtningsmaade er
behersket af denne Aand, sees baade af hvad han bekæmper,
af hvad han erklærer for utilgængeligt for Erkjendelsen, og
af hvad han tager til Udgangspunkt.

Fra Theologien, siger han, har Videnskaben arvet den
Vane at betragte Sjælen som et ulegemligt Væsen, der ikke
kunde opfattes gjennem Sandserne. Derved opstod „den
metafysiske Sjælens Filosofi“, som ikke tog tilbørligt Hensyn til
de fysiske Betingelser, under hvilke Sjælelivet udvikler sig.
Man begyndte der, hvor man skulde ende, idet man fremstillede
Sjælens Filosofi efter Selvbevidsthedens tvivlsomme Aabenbarelser,
medens det dog netop var disse Aabenbarelser, der
skulde forklares. Man undervurderede og foragtede Legemet,
betragtede det som Noget, man maatte skamme sig ved. Derfor
gaaer Maudsleys Anstrengelser ud paa at frigjøre Problemerne
fra den Taage, hvori Metafysiken skal have indhyllet dem.

I Modsætning hertil indskærpes, at Spørgsmaalet om
Sjælens Natur er et saadant, at Videnskaben aldeles ikke kan
tage fat paa det, men at vi ere og rimeligviis altid ville vedblive
at være fuldstændig uvidende i denne Henseende. Vi
kunne kun svagt famle efter Lovene for dens Virksomhed.

Paa hvilken Maade kunne da dens Love findes? Naar
man ikke længer betragter Sjælen som et usandseligt Væsen,
er der Intet til Hinder for at anvende Induktionsmethoden
ved Undersøgelsen af Sjælen ligesaavel som ved
alle andre Naturfænomener. Vi begynde da med de simpleste
Kjendsgjerninger, nemlig de legemlige Betingelser og Udtryk
for Sjælens Virksomhed. Enhver, som vil berige vor Kundskab
om Sjælelivet, maa i vore Dage først og fremmest udforske
de legemlige Betingelser for dets Fremtræden i sund
og syg Tilstand. Hvad vi saa vinde ved Betragtningen af
de simplere Kjendsgjerninger, benytte vi til Forklaring af de
mere indviklede og vanskelige. De høiere Bevidsthedsfunktioner
% --- p. 311 ---
\opage{311}udvikles af de lavere, ligesom en eiendommelig og indviklet
Struktur er en Udvikling af en simplere og mere almindelig.
Af det Konkrete udlede vi det Abstrakte, af det Usammensatte
det Sammensatte. Bevidsthedslivet er Sum eller Facit af
en Strøm af Kræfter, til hvilken ethvert legemligt Organ
sender sine Bidrag, hvad enten disse nu bemærkes eller ikke.
Summen maa opløses i sine Addender, Facitet i sine Faktorer
for at kunne forstaaes.

Man seer, at Maudsley deler Comtes Opfattelse af
Videnskabens Historie og Karakteer. Den theologiske, den
metafysiske, den positive Periode --- denne Tredeling træder
klart frem i det nys Anførte. Og som Comte erklærede
Tingenes Væsen i og for sig, eller, som han for det Meste
udtrykker sig, deres Aarsager, for absolut uerkjendelige, saaledes
stiller Maudsley Spørgsmaalet om Sjælens Natur
udenfor Videnskaben. Hvad der da bliver tilbage som
Videnskabens Opgave, er at undersøge Betingelserne, Lovene
for Tingenes indbyrdes Forhold. Maudsley er endelig ogsaa
enig med Comte i den almindelige Sætning, at det Høiere
skal udledes af det Lavere, det Elementære. Men den Maade,
hvorpaa Maudsley har opfattet Positivismens Principer, og
den klassiske Fremstilling af Exempler, hvormed han oplyser
dem, giver hans Foredrag en særegen Interesse. --- Vi dvæle
først ved de almindelige Principer, førend vi gaae over til
Enkelthederne.

Polemiken mod „Metafysiken“ holdes, ligesom hos Comte,
temmelig ubestemt, og det er derfor ikke let at gaae nærmere
ind paa den. Paa den ene Side synes det, som om der
ved Metafysik skulde betegnes enhversomhelst Besvarelse af
Spørgsmaalet om Tingenes Væsen og Natur, altsaa her
om Sjælens Natur. Men paa den anden Side betegner
Maudsley med hiint Navn særlig den spiritualistiske Opfattelse,
der betragter Sjælen som et immaterielt Væsen, der i
sig selv er uafhængigt af de legemlige Betingelser. Men
selv da er det tvivlsomt, hvorvidt hans Polemik virkelig
rammer. Netop Spiritualisterne have været meest opmærk%
% --- p. 312 ---
\opage{312}somme paa de legemlige Betingelsers Betydning; jo høiere
de hævede Sjælens Væsen over den fysiske Verden, desto
stærkere maatte den Indvirkning, de legemlige Betingelser
øvede paa den, blive. I dette Punkt vilde Spiritualister og
Materialister prægtigt kunne enes. Platon, denne Spiritualismens
Fader, lærer udtrykkelig, at det er Legemets Beskaffenhed,
der frembringer Sjælens Sygdomme, at „ingen Harmoni og
Disharmoni har større Indflydelse paa Sundhed og Sygdom,
Dyd og Last, end den mellem Sjælen og Legemet“; --- ja, han
% as printed (p. 312): no comma after „Materialister“ at the line end
gaaer saa vidt som Nogen af de yderligstgaaende Materialister
naar han siger, at Ingen er slet med sin Villie, men at de „Slette
% footnote *) on p. 312
% in the footnote a small raised mark after „Overs.“ (a speck or displaced spacing material, not a letter) is not transcribed; the dash in „138--141“ has a small tick below it (type damage)
blive slette ved en sygelig Legemsbeskaffenhed og ved en slet Opdragelse“\footnote{Timæus. (Heises Overs. S. 138--141).}. Og hvad det særlige Spørgsmaal om Afsindighed
angaaer, som er det, hvortil Maudsley knytter sin Undersøgelse
af Forholdet mellem Sjæl og Legeme, saa have
i den nyeste Tid netop spiritualistiske Tænkere (som f. Ex.
Franskmanden Albert Lemoine) hævdet, at den nødvendigviis
maa have en legemlig Aarsag, fordi de nægte, at Sjælen i
sig selv kan være syg. --- Skal Polemiken mod Metafysiken
derimod gjælde ethvert Forsøg paa en Løsning af Problemet
om Sjælens Natur, gjælder den ligesaavel Materialismen
som Spiritualismen og etablerer en empirisk
Skepticisme. Vi have ovenfor seet, at dette skeptiske Element
var karakteristisk for Positivismens Udgangspunkt.

Men Maudsley fastholder ligesaalidt som Auguste Comte
denne Skepsis. Ved den Maade, hvorpaa han anvender
Induktionsmethoden, naaer han til en virkelig Besvarelse af
det Problem, han først afviste, „tager virkelig fat paa“ Spørgsmaalet
om Sjælens Natur --- og bliver forsaavidt selv
Metafysiker. Der er ingen Tvivl om, at han har fuldkommen
Ret i at anvende Induktionsmethoden ved Undersøgelsen
af Sjælelivet; men det kommer an paa, hvorledes Forudsætningerne
for Methodens Anvendelse oprindelig stilles: en Feil
her maa nødvendigviis, trods den skarpsindigste og nøiagtigste
% --- p. 313 ---
\opage{313}Gjennemførelse af Methoden, frembringe en Feil i Resultatet.
Maudsley begaaer nu fra Begyndelsen den Feil at betragte
de elementære Betingelser for Sjælelivet og dettes
rudimentære Former som det Konkrete og Simple, hvoraf
Bevidstheden i sine forskjellige Former fremgaae ved Sammensætning;
det høiere Sjæleliv betegnes da i Forhold til det
elementære som en Abstraktion. Og dog bliver det af
Maudsleys egen Fremstilling klart, at det er det omvendte
Forhold, der er det sande. „Det bliver nødvendigt,“ siger
han, „at udskille de lavere og mere elementære af disse
Funktioner og af det Konkrete saa godt som muligt at udlede,
hvorledes Abstraktionen skal forstaaes.“ Men hvad
er denne Udskillen Andet end en Abstraheren? Den enkelte
Funktion drages ud af den hele Sammenhæng, i hvilken
den har sin eiendommelige Tilværelse. Det Elementære, der
udskilles paa denne Maade, bliver abstrakt. Hvad Maudsley
kalder „Abstraktionen“, er den totale Sammenhæng, i hvilken
alle enkelte Funktioner virke, og som omfatter og leder dem,
--- en Sammenhæng, der paa sit høieste Udviklingstrin er
Bevidstheden selv, Sjælen, forsaavidt den, for at bruge
Maudsleys eget Udtryk, „samler og sammenfatter det legemlige
Liv.“ Denne Samlen og Sammenfatten forudsætter en
særegen Aktivitet og godtgjør, at Sjælen ikke blot er en
Abstraktion eller en Sum af Kræfter. Gjennem den totale
Sammenhæng, der ligesaavel betinger de elementære Funktioner,
som disse betinge hiin, godtgjør Sjælelivet sig som
et i sig Konkret og Realt. Og i Tænkningen træder Sjælens
aktive Eenhed endnu klarere frem. Enhver Bevidsthedsakt er
i sig selv en enkelt, udelelig Akt. --- Ved at bortvende Opmærksomheden
fra denne Sjælens og Bevidsthedens reale
Eenhed har Maudsley gjort sig Opgaven utilbørlig let,
ligesaa let, som han begyndte med at gjøre den vanskelig, da
han erklærede Sjælens Natur for aldeles ubegribelig. Og
ikke blot dette. Han er ved den Vending, han har givet
Undersøgelsen, kommet over i den materialistiske Opfattelse
af Sjælelivet. Det er jo det Eiendommelige for denne at
% --- p. 314 ---
\opage{314}sætte det Sjælelige som det blot Afledede (Sum, Facit,
Abstraktion) og Betingede i Forhold til det Legemlige og
ikke at indrømme hiint et eiendommeligt Væsen for sig.

Det er, som sagt, ikke Induktionsmethoden selv, der
fører til dette Resultat, men de Forudsætninger, der lægges til
Grund ved dens Anvendelse, og i hvilke der indsmugles
Bestemmelser af dybt indgribende Betydning. Maudsley
henleder Opmærksomheden paa „de smaa Tings Rige“, de
uendelig fine Nuancer og Bevægelser, som let undgaae Forskerens
Blik, og som dog kunne være af stor Betydning. Men
i logisk Henseende har han selv overseet de „smaa Ting“,
der, skjulte i Forudsætningerne, fremtræde i Konsekvenserne
som meget store Ting. Han begaaer ved Begyndelsen en
logisk, eller om man vil en metafysisk Feiltagelse, som øver
sin Indflydelse paa den hele Udtydning af de fysiologiske
Resultater, men som disse ere uden Skyld i. De fysiologiske
Undersøgelser miste ikke deres Værd, fordi Forudsætningerne
for den afsluttende Erkjendelse, --- en Erkjendelse, der gaaer
ud over Fysiologiens Grændse, --- ere falske. --- Vi skulle
nu fra de mere almindelige Betragtninger gaae til det Enkelte
for ogsaa her at undersøge Forholdet mellem den
positive Forsken og de Resultater, til hvilke den mener at
være naaet.

Det er fra Pflügers og Flourens's bekjendte Forsøg,
Maudsley tager sit Udgangspunkt, og det er om Udtydningen
af disse Forsøg, det gjælder. En Frø, der har
% the d of „det“ (p. 314) is poorly inked and looks like a broken letter; read „det“
mistet sit Hoved, trækker det Been bort, som man kniber, og
hopper bort, hvis Pirringen vedvarer; skjærer man Foden af
paa det Been, med hvilket den afgnider en ætsende Syre,
tager den efter flere forgjæves Forsøg det andet Been
til Hjælp. En Due, af hvilken man borttager den store
Hjernes Halvkugler, taber saa at sige sit Initiativ, men flyver
dog, naar man kaster den i Luften, lukker Øinene for Lyset
og synker den Føde, man skyder ind i Munden paa den.
Ere nu disse saakaldte Reflexbevægelser eller automatiske
% printer's error (p. 314): „ellér“ -- a Fraktur é from the foul case for the second e of „eller“ (second reading; confirmed at 600 dpi in the caller); transcribed as printed
Handlinger sjælelige ellér blot legemlige? At tilskrive Sjælen
% --- p. 315 ---
\opage{315}dem vilde efter Maudsleys Mening være et Slags Anthropomorfisme:
det vilde „smage af den gamle, daarlige Retning,
som har gjort saa megen Skade i Filosofien, og som
gaaer ud paa at forklare Kjendsgjerningerne i Naturen ved,
hvad vi føle i vor egen Sjæl.“ Derimod udleder han af
hine Forsøg, at Handlinger med bestemt Øiemed kunne være
aldeles ubevidste: „Dyret har i sin Rygmarv faaet nedlagt
Evne til at foretage saadanne Bevægelser til Selvopholdelse,
hvilke ere det givne som en Deel af dets Natur, og uden
hvilke det knap kunde leve en Dag“. Ogsaa hos Mennesket
viser der sig noget Lignende. Ikke blot er Nysen og Hosten
en Reflexbevægelse, en Bevægelse, der uafhængig af Villien
tjener til at bortskaffe det for Organismen Skadelige; men
selv de egenlige Villiesakter forudsætte til deres Udførelse
automatiske Handlinger ud fra Rygmarvens Bevægelsescentra;
Bevægelsens Hvorledes hører jo ikke ind under den
bevidste Villie. Og naar Mennesket gjentagne Gange foretager
sig Noget, da opdrages Rygmarven ligesom til at
handle paa egen Haand, saa at Handlingen foretages aldeles
instinktmæssig og automatisk, medens Bevidstheden selv sysler
med ganske andre Ting; saaledes mange daglige Handlinger
(Paaklædning og Afklædning o. Lign.). Det er ifølge
Maudsley ikke nogen ringe Deel, der ved Læren om Reflexbevægelserne
trækkes fra de Fænomener, vi pleie at betragte
som sjælelige.

Men Slutningen synes her ikke at være ligesaa nøiagtig
som de Experimenter, fra hvilke den drages. Saalænge
den ene Mulighed staaer ved Siden af den anden, kan man
dog ikke sige, at Besvarelsen er positiv. Cuvier og Johannes
Müller droge med fuld Ret den modsatte Slutning
% the r of „efter“ (p. 315) has its foot broken off below the stem (a detached stroke at the baseline), not the descender of an x
af Flourens's Forsøg, nemlig at Dyret efter at have mistet
den store Hjernes Halvkugler vel blev sløvt og dorsk, men
dog viste tydelige Tegn paa virkelig Fornemmelse. Hvorledes
vil Maudsley kunne modbevise denne Fortolkning (som
han mærkværdig nok ikke engang omtaler), da den har Tegnene
for sig? Tilmed da han ikke vilde kunne undgaae at
% --- p. 316 ---
% p. 316: „(sic!)“ and „(sic)“ in antiqua
\opage{316}komme ind paa noget Lignende, naar man vilde afæske ham
en nærmere Forklaring af de mystiske Udtryk, han vil sætte
istedetfor Fornemmelse og Villie: „Evne til Selvopholdelse,
til at føle \textit{(sic!)} og undgaae, hvad der er skadeligt, ligesom
ogsaa til at føle og stræbe \textit{(sic)} efter, hvad der er velgjørende.“

Sagen er, at Maudsley ved den ovenfor omtalte „Udskillen“,
som han ikke saae var en Abstraktion, har sat en
skarp Forskjel mellem det sjælelig-bevidste og det organiske,
ubevidste Liv. Men en saadan Forskjel kan ikke fastholdes
uden mangfoldige Overgangspunkter. Indenfor Fornemmelsen
selv er der en uendelig Skala af Trin, som føre fra det
ubevidste til det bevidste Liv. Dette kan derfor ikke udledes af
hiint, uden ved at forudsættes som Anlæg og real Mulighed
deri. Der er et indre Forhold mellem det sjælelige Princip
og dets Betingelser, saaledes at disse fra første Færd indeslutte
Principet i sig, for Skridt for Skridt at lade det
udvikle sig, idet de udvikle sig selv. Bevidsthedens Grad
kan altsaa vexle; men naar man, som Maudsley selv indrømmer,
har Handlinger for sig, „som i Virkeligheden
bære Præg af forudgaaende Bevidsthed og Villie“, er der
ingen Grund til at mene, at Bevidstheden aldeles skulde
% footnote *) on p. 316
% emphasis in the footnote: „Virchow“ letterspaced (gaps 16, 11, 12, 10, 9 px against 1--6 px in „bruger“) and „bevidste“ letterspaced (gaps 8, 7, 8, 9, 7, 10, 7 px against 1--5 px solid); „p.“ in antiqua; the German title „Ueber das Rückenmark“ in Fraktur
mangle\footnote{\emph{Virchow} bruger til Forklaring af Reflexbevægelserne Begrebet „ubevidst Fornemmelse“. (Ueber das Rückenmark. Berlin 1871 \textit{p.} 25). Herved angives den rette Mellembestemmelse, naar man erindrer, at Fornemmelse aldrig kan staae i fuldstændig Modsætning til Bevidsthed. --- Han gjør iøvrig opmærksom paa, at de Paavirkninger, der fremkalde Reflexbevægelserne, forplantes gjennem Følenerverne, og at der jo desuden ogsaa er Reflexbevægelser, som foraarsages ved \emph{bevidste} Fornemmelser (f. Ex. Sammentrækning af Øinene paa Grund af for blændende Lys).}. De automatiske Handlinger hos Mennesket øve jo
desuden en betydelig Indflydelse paa selve Bevidsthedslivet;
men hvorledes kunde det skee, naar de foregik aldeles udenfor
Bevidstheden? Sjælen kan ikke „samle og sammenfatte“,
hvad der falder udenfor dens Omraade.

% --- p. 317 ---
% p. 317: the page opens a new paragraph (indented); p. 316 ended with a full paragraph and a complete footnote, so nothing is carried over
% emphasis: „Forestillinger“ letterspaced (gaps 8, 8, 8, 7, 13, 7, 11, 10 px against 1--5 px solid)
\opage{317}Maudsley bliver ikke staaende ved de blotte Reflexbevægelser.
Hvad han her har troet at vinde som Resultat,
søger han at anvende paa de høiere sjælelige Funktioner.
Trods Funktionernes Forskjellighed mener han dog, at Egenskaber
og Love ere væsenlig de samme som ved de lavere
Nervecentra: „Indtryk modtages, paa disse Indtryk følger
der en Reaktion, og Virkningerne saavel af de første som af
den sidste ordnes og bevares.“ Indtrykkene ere de fysiologiske
Betingelser for \emph{Forestillinger}. „Fornemmelsen af Forestillingerne
% emphasis: „Følelse“ (first occurrence) letterspaced (gaps 9, 9, 9, 8, 8, 8 px); the second „Følelse“ in the line is solid
er \emph{Følelse}; ved Følelse forstaaer jeg nemlig
Nervesubstansens særlige Modtagelighed for Forestillinger.
% emphasis: „Hukommelse“ (gaps 9, 12, 11, 9, 8, 11, 10, 7, 13 px) and „Villien“ (gaps 9, 9, 10, 7, 8, 7 px) letterspaced
At disse ordnes og bevares er \emph{Hukommelse}, og i den
Reaktion, de fremkalde, er \emph{Villien} givet.“

Her have vi da i kort Begreb, hvad der skal træde i
Stedet for den „metafysiske Sjælens Filosofi“, som skal
have anrettet saa megen Skade. Saa Meget er vist, at disse
Bestemmelser ikke ere metafysiske; en anden Sag er det, om
de derfor have Mindre af det Taagede, der efter Maudsleys
Mening er eiendommelig for Metafysiken.

Det Første, der falder i Øinene ved hine Bestemmelser,
er, at det aldeles ikke siges, hvad en Forestilling er. Der
angives kun, at Indtrykkene ere dens fysiologiske Betingelse,
og dermed betragtes dens Væsen som forklaret. Denne
Overfladiskhed har sin Grund i, at man for at forklare
Forestillingens Væsen nødvendigviis maa forudsætte Bevidstheden;
men det kan Maudsley, der netop vil bevise, at Bevidstheden
fremkommer som Facit eller Sammensætning, ikke.
Han holder sig da til de empiriske Betingelser og anseer
% footnote *) on p. 317
% the footnote is set in antiqua throughout, including its quotation marks; „Stuart Mill“ spaced antiqua (Stuart: gaps 7, 8, 10, 9, 7 px; Mill: 13, 11, 11 px against 2--5 px solid)
dem for en tilstrækkelig Forklaring\footnote{\textit{„It is one question what a thing is, and another what it depends on“.} \emph{\textit{Stuart Mill}}. \textit{System of logic. II p. 7 (5th. ed.)}}. --- Sagen bliver ikke
klarere, fordi det et andet Sted hedder, at „Forestillingerne
tilegnes gjennem Sandserne, blandes, sammenbindes og grupperes
--- og disse organiske Kombinationer ere de fysiologiske
% --- p. 318 ---
\opage{318}Betingelser for de høieste Former af vor aandelige Virksomhed:
Reflekteren, Ræsonneren og Dømmen.“ Da vi
iforveien kun vide saa Meget om Forestillingerne, at de beroe
paa fysiologiske Betingelser, er det ikke let at sige, hvad en
Blanden, Sammenbinden og Grupperen af dem skal betyde.
Dog --- naar vi videre høre, at dette er „organiske Kombinationer“,
der igjen blive „fysiologiske Betingelser“, saa
kunne Forestillingerne selv ikke være Andet end organiske
% emphasis: „beroe“ letterspaced (gaps 10, 8, 10 px)
Processer, \emph{beroe} altsaa ikke blot paa saadanne, da en saadan
Beroen ikke udelukker væsenlig Forskjel. Hermed vilde Bestemmelsen
% the dot before „Forestillinger“ is a speck
af Følelsen stemme: „Nervesubstansens særlige
Modtagelighed for Forestillinger“, --- Forestillingerne maae
være noget Legemligt for at kunne paavirke Nervesubstansen.
Men den samme Vanskelighed kommer her igjen: Ingen kan
dog for Alvor mene, at der af Nervesubstansens Modtagelighed
følger det Mindste til Forklaring af et Fænomen
som Følelsen, der er en indre Tilstand og forudsætter et
Forhold af Væsenet til sig selv.

Overfor disse Vanskeligheder har Maudsley ikke Ret
til at paaberaabe sig Uerkjendeligheden af Sjælens Væsen.
% printer's error (p. 318): the last letter of „tidligere“ is a wrong sort shaped like a 2 (no e-stem, no descender; confirmed by the second reading); the word is transcribed „tidligere“
Thi vi have allerede tidligere seet, at han ikke selv fastholder
denne Uerkjendelighed. Og han gjør det heller ikke i sine
fysiologiske Bestemmelser. Naar det f. Ex. hedder: Følelse
er Nervesubstansens Modtagelighed o. s. v., saa fremtræder
denne Sætning med samme positive Myndighed som hans
øvrige Lærdomme. Derimod følger af hine Vanskeligheder
Nødvendigheden af ved Undersøgelsen af de sjælelige Fænomener
ikke blot at tage sit Udgangspunkt i de fysiologiske
Betingelser, men indenfor det Sjælelige selv, i Bevidstheden.
Og dette er ikke et nyt og fremmed Udgangspunkt. Bevidsthedslivet
med sine eiendommelige Love viser saavel ved
sin sukcessive Udvikling, sine Kriser og Sygdomme som ved
selve sit Formaal, Tilegnelsen og Udviklingen af Tilværelsens
Indhold tilbage til reale, fysiske Betingelser som sit
Grundlag. Det Psykologiske viser i sig selv hen til det
Fysiologiske. Har Fysiologen derimod ikke en klar Bevidst%
% --- p. 319 ---
\opage{319}hed om det Psykologiskes Eiendommelighed, vil Spørgsmaalets
Kjerne undgaae ham.

At denne Bevidsthed mangler hos Maudsley, sees ikke
mindst af hans Lære om Hukommelsen. Denne bestaaer
ifølge ham i en Ordnen og Bevaren af Forestillinger. Men
det er, vedbliver han, en metafysisk Vildfarelse, at Hukommelsen
skulde være eiendommelig for Sjælen. Rygmarven
og de sensoriske Gangliers erhvervede Funktioner kunne ikke
gaae for sig uden Hukommelse; disse Organer vilde blive
idiotiske, naar der ikke var Hukommelse i dem. „I enhver
Nervecelle er der Hukommelse, og, hvad Mere er, den findes
i ethvert organisk Element af Legemet. Koppe- eller Syfilisgiften
sætter sit Mærke paa Legemet for den øvrige Deel
af Livet; vi glemme det maaskee, men det glemmer ikke os.
Den Maade, hvorpaa Arret efter et skaaret Saar i Fingeren
paa et Barn holder sig og voxer med Barnets Legeme,
godtgjør, at de organiske Bestanddele af Fingeren huske den
Forandring, de ere undergaaede. Hukommelsen er en ordnet
Bevarelse af de Virkninger, Indtrykkene have fremkaldt.“
Men hvor komme vi nu hen? Uformærkt tabes af Sigte
den første, psykologiske Bestemmelse af Hukommelsen, og al
Ordnen og Bevaren af Virkninger efter Indtryk gjøres til
Hukommelse. Der tages ikke Hensyn til den væsenlige Forskjel
mellem den Maade, hvorpaa Forestillingerne opbevares
i Bevidstheden, og den, hvorpaa Virkninger bevares i Organismen.
Et Navn, der indridses i Barken af et Træ,
bevares dog ikke af dette saaledes, som det bevares af den,
hvis „kjære, stadige Tanke“ det er. Et af To, enten maa
Maudsley her paa Digterviis gjøre Træet til et besjælet
Væsen med Erindring og Sympathi, eller ogsaa maa han
gjøre Bevidstheden til en reent mekanisk Beholder. Hvad
der frembringer Uklarheden er, at Hukommelsen er utænkelig,
uden at de organiske Elementer fastholde den til den huskede
Forestilling svarende Bevægelse, eller rettere dennes Mulighed.
Men denne organiske Bevægelse er ikke Forestillingen selv,
kun dens Betingelse. I det organiske Element „huskes“ den
% --- p. 320 ---
\opage{320}til Forestillingen svarende Bevægelse paa samme Maade som
Arret i Drengens Finger, Navnet i Træets Bark; i Bevidstheden
huskes Forestillingen som Moment i den hele
Tankeverden, der er Bevidsthedens Indhold, og den huskes
med en Intensitet, der beroer paa dens Betydning i dette
Indhold. Saadanne Ord som Indtryk, Orden, Samlen
faae derfor en forskjellig Betydning efter den Side, hvorfra
man seer Sagen. Men Positivisten, som overseer dette,
kommer ikke blot til at tale fysiologisk, hvor han skulde tale
psykologisk, men ogsaa psykologisk, hvor han skulde tale fysiologisk
--- og gaaer derved tilbage fra det „positive“ Standpunkt
til det „metafysiske“, ja maaskee endog til det „theologiske“,
idet han forklarer Naturfænomenerne efter sit eget
Væsens Analogi.

Ligesaalidt som Hukommelsen er, efter Maudsley, Villien
et eiendommeligt sjæleligt Fænomen. Det er nemlig, siger
han, en Grundeiendommelighed hos alle organiserede Væv,
% the small tick after „der“ (much smaller than a comma) is a speck, not transcribed
at de stræbe efter, hvad der er gavnligt, og skye, hvad
der er skadeligt for dem. De Tilskyndelser, som opstaae af
denne Egenskab ved Vævene, kalde vi Villiesyttringer, naar
de ledes af Fornuften. --- Bortseet fra den allerede tidligere
omtalte Vanskelighed ved at tænke sig, hvad det skal sige, at
et organiseret Væv stræber, skyer og tilskynder, --- er der
det Eiendommelige ved denne Bestemmelse af Villien, at det
Afgjørende indføres som en Betingelse, der ikke forklares.
Hvad der adskiller Villien fra den blotte Drift, er ganske
vist det, at den ledes af Fornuften; men det maa saa først
og fremmest vises, af hvad Art dette Forhold mellem Fornuften
og Drift-Tilskyndelsen er, og hvorledes det er muligt.
Hvorledes kunne Tilskyndelserne paa eengang fremkaldes ved
de organiserede Vævs Egenskaber og „ifølge Reflexionen
ledes af Fornuften“? Dette Spørgsmaal, som psykologisk
er Hovedspørgsmaalet, opkaster Maudsley slet ikke, ligesaalidt
som vi faae at vide, hvad han forstaaer ved Fornuft. Derved
kommer han til at mangle en virkelig Maalestok for
Forholdet mellem Reflexbevægelserne og de egenlige Villies%
% --- p. 321 ---
\opage{321}yttringer; thi naar man ved disse sidste tilsidesætter Fornuft-
% the wider gap between k and e in „ganske“ is a loose sort, not letterspacing (the other gaps are solid)
eller Bevidstheds-Elementet, blive de ganske vist ikke
forskjellige fra hine.

Man kan, hedder det fremdeles, fysiologisk seet kun
% emphasis: „Villien“ letterspaced (gaps 10, 8, 9, 9, 10 px)
tale om Villiesyttringer; \emph{Villien}, som Noget, der er forskjelligt
fra de enkelte Handlinger, er en metafysisk Abstraktion.
--- Om der virkelig er Nogen, der har villet stille
Villien udenfor Handlingerne, maa vist betvivles; Maudsleys
Polemik er her som paa andre Punkter ikke fri for selv
at frembringe den Modstander, han bekæmper. Man vilde
i ethvert Tilfælde kunne rette den samme Polemik mod ham
selv. Den Grundeiendommelighed, han antager hos de
organiserede Væv, bliver ved at anlægge hans Betragtningsmaade
ogsaa kun en „metafysisk Abstraktion“; thi den yttrer
sig kun under bestemte Paavirkningers Indflydelse. Men
hvad forstaaer man ved Villie Andet end den Grundeiendommelighed
ved Sjælelivet, hvortil Villiesyttringerne vise
tilbage? Den positive Forsken, der holder sig til de faktiske
Virkninger for ved dem at lære Tingenes Væsen at kjende,
kan dog, som Maudsley's eget Exempel viser, ikke undgaae
at føre Virkningerne tilbage til Væsenet som Muligheder,
der indeholdes deri. Dette er maaskee metafysisk; men det
er saa en Metafysik, til hvilken Fysiken selv fører os. Vil
man f. Ex. sige, at vi ikke kjende Elektricitetens Væsen,
maa dette begrændses derhen, at vi kjende Elektriciteten,
forsaavidt vi vide, hvorledes den virker under bestemte Forhold.
% the faint low mark between „et“ and „Bindeled“ is inked spacing material, not a hyphen
Vi vide jo saaledes, at den elektriske Strøm danner
et Bindeled mellem de forskjellige Naturvirksomheder, idet
man ved den kan frembringe Lys og Varme, magnetiske og
kemiske Virksomheder, og af disse omvendt hiin; og man
kan maale den elektriske Strøm ved dens kemiske eller magnetiske
% footnote *) on p. 321
% emphasis in the footnote: „Holten“ letterspaced (gaps 11, 9, 9, 6, 7 px); the footnote is in Fraktur (digits roman)
Virkning\footnote{C. V. \emph{Holten}: Om Virksomhedernes Omdannelse i Naturen. 1868. (Universitetsprogram). S. 19.}. Skulde dette ikke være et Bidrag til
Erkjendelse af Elektricitetens Natur? Søger man Væsenet
% p. 321: the „21“ at the foot is the sheet signature, not transcribed
% --- p. 322 ---
\opage{322}eller Naturen udenfor Yttringerne, gjør man sig skyldig
i en falsk Metafysik. Vi have jo oplevet det Særsyn, at en
Tænker som Stuart Mill har skrevet en Bog om Kvindens
Stilling i Samfundet, som indeholder den sindrigste og meest
træffende Karakteristik af Kvinden --- men hvori han dog
fra først til sidst fastholder, at vi Intet vide om Kvindens
Natur!

% p. 322: a short centred rule (c. 1.3 cm; ink as dark as the text, not show-through) and the numeral „4.“ centred in roman type
\begin{center}\rule{1.8cm}{0.4pt}\end{center}

\secnum{4.}

Deels som Middel til nærmere at oplyse, hvad han
har udviklet om Forholdet mellem Sjæl og Legeme i Almindelighed,
deels som en særegen Anvendelse deraf, gaaer
Maudsley ind paa en Undersøgelse af Sindssygdommenes
Væsen og deres Forhold til andre Nervelidelser. Disse
Undersøgelser ere de interessanteste og meest belærende i
hans Foredrag; det er tillige de, hvor han meest undgaaer
en uheldig Overskriden af sin Videnskabs og dermed ogsaa
af sin Kompetences Grændse. Naar vi dog ogsaa her
forsøge at anlægge en Kritik, maa det endnu mere end
paa andre Punkter fastholdes, at denne ene gaaer ud fra et
psykologisk Synspunkt og altsaa ikke kan gjøre noget Skaar
i den Betydning, Maudsleys Undersøgelser kunne have i
medicinsk Henseende.

Den sædvanlige Opfattelse og Inddeling af Sindssygdommene
tilfredsstiller ikke Maudsley, fordi den udelukkende
gaaer ud fra et psykologisk Standpunkt. Den sande videnskabelige
Aand, „Positivitetens Aand“, tilfredsstilles efter
hans Mening først, naar man fastholder, at disse
Lidelser beroe paa de samme Vilkaar som andre Sygdomme.
Da Sjælelivet er saa nøie knytet til og afhængig
af de legemlige Funktioner, kan Sindssygdom
heller ikke være et reent sjæleligt Fænomen; det er ikke en
Sygdom hos et metafysisk Begreb, men en nervøs Sygdom
med overveiende sjælelige Symptomer. Denne Anskuelse er
naturligviis ikke ny; hos os har saaledes allerede F. G.
% --- p. 323 ---
\opage{323}Howitz indskærpet den i sin Afhandling „om Afsindighed og
Tilregnelse“ (1824), hvor Afsindighed bestemmes som „en
Indskrænkning af Fornuften eller Fornuftens Brug formedelst
en Sygdom i de materielle Organer for Sjælens Virksomhed“,
som „Hjernesyge“, og hvor der derfor lægges Vægt
paa den dobbelte (sjælelige og legemlige) Undersøgelse og
Bestemmelse af den. Men det Eiendommelige ved Maudsleys
Undersøgelse er den nærmere Begrundelse og de videre
Konsekvenser.

Man kjender aldeles ikke de Bevægelser i Hjernedelene, som
ere de fysiske Betingelser for den sjælelige Virksomhed, eller
de kemiske Forandringer, som ledsage dem, ligesaalidt som
man kan paavise nogen Forskjel mellem den trætte, afkræftede
og den oplivede, forfriskede Hjernes Nerve-Elementer.
Undersøgelsesmidlerne slaae ikke til overfor Forandringer, der
ere finere end de mikroskopiske. „Vi kunde næsten ligesaa
godt give os til at studere Anatomien af en Myg med et
Teleskop“. Man har saaledes fundet afdøde Sindssyges
Hjerner ganske normale. Men man maa efter Maudsley
ikke slutte heraf, at Forandringen ikke er der; den kan godt
være til, om den end hører til „de uendelig smaa Tings
Rige“, som er utilgængeligt for vore Sandser. I Kraft af
den nøie Forbindelse mellem Sjæl og Legeme og af de
% emphasis: „maa“ letterspaced (gaps 9, 9 px against 3--6 px in „Sindssygdom“ and „der“ on the same line)
legemlige Symptomer, der ledsage Sindssygdom, \emph{maa} der
antages en Sygdom i Hjernen.

Denne Slutnings Berettigelse staaer dog ikke ganske
uanfægtet. Man maa for det Første undre sig over, at
Maudsley ikke tager noget Hensyn til de Tilfælde, som ere
konstaterede, hvor begge den store Hjernes forreste Lapper
vare syge og ødelagte i en betydelig Grad, og hvor der dog
% footnote *) on p. 323
% emphasis in the footnote: „Longet“ (gaps 9, 8, 9, 9, 7 px) and „Lange“ (gaps 8, 7, 9, 11 px) letterspaced; the German title is in Fraktur
% as printed in the footnote: „Geschichte der Materialismus“ (der for des); not corrected
ikke blev iagttaget den mindste Forstyrrelse i Intelligensen\footnote{\emph{Longet} anfører to saadanne Tilfælde. Se \emph{Lange}: Geschichte der Materialismus. S. 431.}.
Og dernæst vil man mod Maudsleys Tankegang kunne
stille en anden, som det synes, mere videnskabelig-besindig, af
% --- p. 324 ---
\opage{324}den franske Fysiolog Leuret. Ganske vist, siger denne, har
jeg ingen Ret til at slutte, at der ingen Forandring er,
fordi jeg ikke kan see nogen; men med samme Forsigtighed
% emphasis: „er“ letterspaced (gap 8 px against 2--3 px in „der“ before it)
maa jeg vogte mig for at slutte, at der \emph{er} nogen. Naar
en Afsindigs Hjerne synes at være sund, saa slutter jeg ikke,
at denne Hjerne er syg, men jeg forbliver i Tvivlen, indtil
Sandheden er beviist. Og dersom de Tilfælde, hvor Hjernen
synes at være sund, netop ere de, hvor Afsindigheden ikke
ledsages af legemlige Symptomer, men er en Vildelse i
Forstanden og Lidenskaberne, medens de Tilfælde, hvor
Hjernen er forandret, ere de, hvor der har fundet Lamhed,
Sløvhed og Søvnløshed Sted, saa tilskriver jeg Hjernebeskadigelsen
disse Omstændigheder, og Aarsagen til Sindsforvildelsen
% footnote *) on p. 324
% the footnote is set in antiqua
vedbliver at være ubekjendt\footnote{\textit{Revue des deux mondes. 15 juillet 1865. p. 417.}}.

Saa Meget synes i det Mindste at fremgaae heraf, at det
er forbundet med de største Misligheder at drage for hurtige
Slutninger paa dette Omraade. Men i ethvert Tilfælde
--- Rigtigheden af, at der altid er en legemlig Aarsag indrømmet,
--- har man saa derved vundet en nærmere Indsigt i
Sindssygdommens Væsen, og maa dette ikke vedblivende
nærmest opfattes fra den psykologiske Side? Maudsley bemærker
selv, at de sjælelige Symptomer ere de overveiende.
Men hvorledes kan man paavise en indre Forbindelse mellem
dem og de legemlige Symptomer? Man har med Rette
sagt, at naar et Uhr slaaer forkeert, fordi et Hjul er i
Stykker, saa følger deraf ikke, at det var dette Hjul, der
slog; --- vi kjende ikke et Lokomotivs Natur eller Virksomhed,
fordi vi vide, at det ikke kan gjøre Tjeneste, naar det er i
Stykker. Forbindelsen mellem de sjælelige og legemlige
Symptomer bliver en reent empirisk, hvorimod vi indenfor
Sjælelivet kunne forfølge Afsindigheden Skridt for Skridt
fra de første svage Tegn til den dæmonisk overvældende
Magt, hvormed den optræder paa sit Høidepunkt. Her
% --- p. 325 ---
% p. 325: begins with a whole word („kommer“; p. 324 ends „… Høidepunkt. Her“); no footnote continued from p. 324 at the foot
\opage{325}kommer Ideeassociationens, Tankens og Karakterens Love os
tilgode --- men Studiet af disse Love er et psykologisk Studium.
Sindssygelægen maa derfor fremfor Alt være praktisk
Psykolog.

Afsindighedens legemlige Aarsag søger Maudsley ikke i
en Paavirkning udenfra, ved aandelig Rystelse eller legemlig
Beskadigelse; en saadan Paavirkning er kun Anledning.
Sindssygdommens egenlige Rod ligger dybere, i et medfødt
Anlæg, i Organisationen af Individets Nervesystem. Naar
to Mennesker udsættes for samme aandelige Rystelse, og den
Ene af dem bliver afsindig, medens den Anden kun for et
Øieblik mister sin Ligevægt, kan Grunden dertil kun
være, at der hos hiin er en Faktor med i Spillet, som ikke
findes hos denne, --- en eller anden medfødt Nervesygdom.
Vilde man mene, at denne dog ikke behøvede at være medfødt,
men kunde have dannet sig under Indflydelsen af Livsforholdene,
afviser Maudsley dette paa følgende Maade:
„Af alle de feilagtige Begreber om Sjælen, som Metafysiken
har avlet eller været medskyldig i, er der maaskee intet, som
er mere falsk end det, der stiltiende antager eller udtrykkelig
erklærer, at alle Mennesker ere fødte med lige gode Betingelser
for Udviklingen af et Sjæleliv, medens Omstændighederne
og Opdragelsen betinge Forskjellighederne i den senere
Udvikling“. Den stakkels Metafysik, der nu foruden sine
egne Synder ogsaa maa tage Materialismens Forseelser paa
sine Skuldre og saaledes her anklages for det aldeles Modsatte
af, hvad der ovenfor fremførtes imod den! Dog, den
have vi ikke her at forsvare. Vi skulle derimod see, hvad
Maudsley opstiller i Modsætning til det, han lægger den
til Last. „Opdragelse kan ganske vist udrette Meget, og
Omstændighederne i Livet kunne ogsaa udrette Meget; men
vi maae ikke glemme, at det Grundlag, hvorpaa Opdragelsens
Udbytte maa hvile, ikke tilegnes, men tages i Arv. Ingen
kan undgaae at beherskes af sin Organisme, Ingen kan
undflye den Skjæbne, han er født med“. Det Medfødte er
altsaa nedarvet, og det nedarvede Fond af Kræfter, Mennesket
% --- p. 326 ---
\opage{326}møder med til Livets Kamp, betinger det Held, hvormed
han udfægter denne. Sjælens Liv er som alt Liv stillet
mellem Kræfter, der altid søge at tilintetgjøre det, og mod
hvilke det derfor stedse maa være paa Krigsfod. Den, der
kan beherske Omstændighederne, kommer frem; den Svage
beherskes af dem og synker sammen --- og en af de Maader,
hvorpaa hans Ulykker kunne vise sig, er netop som Sindssygdom.
„Saaledes gaaer det til, at sjælelige Prøvelser,
som tjene til at styrke en kraftig Natur, nedbryde en svag Natur,
som ikke kan vise tilbørlig Modstand; den Omstændighed, at
en moralsk Aarsag kan fremkalde en Sindssygdom, røber en
Sammensværgelse indenfra med de ugunstige ydre Omstændigheder.“

Vi arve som en naturlig Gave, hvad vore Forfædre
have erhvervet, baade det Gode og det Onde. Den moralske
Sands opstaaer ved, at det, som den ene Slægt opnaaer
af bevidst Erkjendelse af det, der er Menneskeheden gavnligt
eller skadeligt, i Kraft af Nervesystemets Organisation (jfr.
ovenfor om Hukommelsen) nedarves til den næste Slægt
som „en naturlig eller instinktmæssig Modtagelighed, der tilhører
hans Hjerne, naar denne har sin rette Bygning“. Det
enkeltviis Erhvervede sammensættes saa nøie, at det tilsidst
synes at udgjøre een enkelt, udelelig Evne. Ved at forfølge
Menneskehedens Historie tilbage vilde man tilsidst, ved Udgangspunktet,
finde „en Tilstand, som, hvad enten man vilde
kalde den abeagtig eller menneskelig, savnede al videre moralsk
Følelse.“ Det fysiologiske Beviis herfor skal ligge i den
menneskelige Hjernes fremskridende Udvikling; den Deel af
Vævet, som er kommen til, repræsenterer Slægtens Erfaringer
og Erindringer, som ere indlemmede i os. Men
det er kun, naar de normale Forudsætninger ere tilstede, at
den moralske Sands nedarves. Er Hjernens Bygning ikke
den rette, har Individet „det sindssyge Temperament“, saa
er det netop hiin Sands, Kulturens høieste Blomst, der først
mistes. Vedvarer dette Anormale gjennem flere Slægtled,
fremkommer Idiotismen. Den, der er født Idiot, har
% --- p. 327 ---
% as printed (p. 327): „theoride“ (for therioide/theroide, 'animal-like'); not corrected
\opage{327}mistet sin Fødselsret som Menneske og synker tilbage paa
det dyriske Trin. Af særlig Interesse ere her de „theoride
Idioter“, hos hvem ikke blot dyriske Instinkter, men ogsaa et
dyrelignende Udseende forekommer --- Træk, der ere ligesom
Ekko fra en forsvunden Tid, Vidnesbyrd om et Slægtskab,
Mennesket næsten er voxet fra og er for stolt til at vedkjende
sig. I saadanne forulykkede og vanslægtede Individer mistes
altsaa det vundne Resultat, og dets oprindelige, ubearbeidede
Element fremtræder igjen. Saaledes bekræfter Sindssygdommenes
Psykologi Darwins Anskuelse, idet den viser tilbage
til den som til den eneste grundige Forklaring.

Det er med en vis Nødvendighed, at Maudsley slutter
sine Undersøgelser med at henvise til Darwinismen, denne
den moderne Naturvidenskabs dristige, afsluttende Anskuelse.
Men det er ikke sagt, at denne Henviisning er Et med en
Forklaring. Darwinismen har en vis Ubestemthed i sit
Væsen, kan opfattes fra en dobbelt Side. At alle organiske
Former have udviklet sig af een Urform, kan nemlig enten
betyde, at alle hine Former kun ere Udtryk for den forskjellige
Maade, hvorpaa Urformen mekanisk paavirkedes og
modificeredes af de ydre Betingelser, eller at de oprindelig
indeholdtes som Anlæg og Muligheder i Urformen, saa at
denne i sig som i en Spire omfattede den hele Verden af
levende Væsener. I sidste Tilfælde bliver Darwinismen ingen
særlig positivistisk eller materialistisk Lære. Thi de høiere
Former ere saa anlagte fra først af og have som Anlæg en
ideel Overlegenhed over Betingelserne, der kun blive Midlerne
til deres Virkeliggjørelse. Og lægger man Mærke til, at den
egenlige Hovedfaktor hos Darwin ikke bliver de ydre Betingelser,
men hvad han kalder Kampen for Livet, synes det
rimeligt, at denne sidste Opfattelse af Darwinismen maa
være den rette. Til Kamp for Livet fordres der et virkeligt
Livsprincip, hvor Energien kan koncentreres, ud fra hvilken
Kampen kan føres. Der maa være Noget, for hvis Bevarelse
% emphasis: „oprindeligt“ letterspaced (gaps 12, 13, 12, 10, 9, 9, 8, 12 px against 1--5 px solid)
og Udvikling der kæmpes; uden et saadant \emph{oprindeligt} Gode
vilde Kampen ikke have ført til Noget. Af Intet kommer Intet.

% --- p. 328 ---
\opage{328}Men Maudsley gaaer netop ud fra et oprindeligt Intet,
fra et Nulpunkt. Slægtens Udvikling skal være begyndt paa
et Stadium, hvor al videre moralsk Følelse manglede. Først
efterhaanden opsamles de Brudstykker, hvoraf den moralske
Sands sammensættes saaledes, at den tilsidst kan tage sig
ud som en udelelig Evne. Men selve Tilegnelsen og Bearbeidelsen
af disse Brudstykker forudsætter en Aktivitet, en
Arbeiden, en „Samlen og Sammenfatten“, --- en Virken
kort sagt, hvis Mulighed vi ikke kunne indsee, naar der ikke
er et oprindeligt Virkende. I psykologisk Henseende gaaer
Maudsley ud fra den Anskuelse, som de engelske Filosofer
% p. 328: the mark after „Nutiden“ is a speck, not a footnote mark or punctuation
i Nutiden (Stuart Mill, Herbert Spencer) have udviklet
som en Arv fra deres store Forgænger i forrige Aarhundrede,
David Hume, --- den Anskuelse nemlig, ifølge hvilken Bevidsthedslivets
Love og hele Udvikling hviler paa Ideeassociationerne,
efter hvilke de empirisk optagne Forestillinger gruppere
sig. De engang grupperede og optagne Forestillinger
blive efterhaanden en varig Deel af den sjælelige Konstitution
og gaae i Arv fra den ene Slægt til den anden, saa at de
for følgende Generationer fremstille sig som i sig gyldige Sandheder
eller gjøre sig gjældende som umiddelbare Instinkter.
% antiqua: „growth of intelligence“ (the parentheses are roman)
Men denne Lære om Forstandens Væxt (\textit{growth of intelligence})
overseer, at den væsenlige, lovbestemte Sammenhæng,
som er eiendommelig for den tænkende Bevidstheds Indhold,
ikke kan være fremkommen ved en blot Ophoben af Forestillinger
eller ved Stadigheden af Indtrykkene udenfra. Den
indre Lovbestemthed er Frugten af et Tankearbeide, der nødvendigviis
forudsætter en arbeidende Tanke, hvis Væsen er
Et med den almene Fornuft i Tilværelsen. Uden dette, det
sidste Princip for al menneskelig Erkjendelse, vilde enhver
Viden være umulig. Og derhos, da de umiddelbare Sandheder
i Tidens Løb blive Gjenstand for Kritik, da Bevidstheden
altsaa kan vende sig mod sit eget Værk for at tilintetgjøre
det, maa der være en indre Eenhed i Selvet, hvis det ikke
under en saadan Kritik skal gaae til Grunde, men skal kunne
arbeide sig frem til nye Udtryk for Sandheden.

% --- p. 329 ---
\opage{329}Den moralske Sands er ikke Et med det, som paa bestemt
Tid og Sted udgjør dens Indhold. Maudsley gjør
sig her selv skyldig i, hvad han bebreider „Metafysiken“, at
tilskrive det Ydre for Meget. Det egenlig Moralske er
væsenlig et Forhold mellem Selvbevidstheden paa den ene
Side, dens Indhold og Betingelser paa den anden Side,
men det bestaaer ikke i dette eller hiint bestemte Indhold. Australnegeren
har derfor ligesaavel moralsk Sands som den i
% p. 329: the n of „Men“ (before „den positivistiske“) is a broken sort, reading like „Mer:“; the word is Men
Intelligens meest udviklede Europæer. Men den positivistiske
og materialistiske Betragtningsmaade gjør sig overfor
saadanne Spørgsmaal skyldig i megen Overfladiskhed. Naar
man f. Ex. hos et vildt Folkeslag ikke finder Tro paa
bestemte Guder eller nogen egenlig Kultus, saa uddrages
strax den Slutning, at dette Folk mangler Religiøsitet, uden
at man undersøger, om det Religiøse ikke hos dette Folk udtrykker
sig paa anden Maade, om det ikke f. Ex. gaaer i
Et med Pietetsforholdet til Høvdinger og Fædre, der altsaa
blive som et Slags jordiske Guder. Eller man frakjender
visse vilde Folk Erkjendelsesevne, fordi man i deres Sprog
ikke finder almindelige, videnskabelige Sætninger udtalte.

% as printed (p. 329): the opening quotation mark before „Arvelighed“ is a single low comma-shaped mark (‚), closed by a double “; transcribed as printed, not corrected
Der skjuler sig fremdeles under Udtrykket ‚Arvelighed“
en Tvetydighed. Et Bevidsthedens Indhold kan ikke ligefrem
arves; det maa tilegnes paany af hver Slægt. Her
kan kun være Tale om umiddelbar Arv, forsaavidt den opvoxende
Generation gjennem Opdragelsen uden Kritik optager,
hvad den foregaaende overgiver den; selv efterat Kritiken er
vaagnet, kan det være vanskeligt at frigjøre sig fra de umiddelbart
modtagne Indtryk. Hvad der derimod ligefrem medfødes,
som en Deel af Individets oprindelige Konstitution,
det kan kun være mere eller mindre gunstige Anlæg og
Drifter. Som det Væsenlige ved den uddannede moralske
Sands er et ethisk Hvorledes, saaledes er det Væsenlige ved
den medfødte Mulighed et organisk-psykisk Hvorledes. Her
er det egenlig først, at Darwins Theori faaer sin Anvendelse,
--- en Anvendelse, som ligger klart for Dagen, og
som allerede Platon har gjort i sin „Stat“, hvor han paa
% --- p. 330 ---
% footnote *) on p. 330
\opage{330}en Maade fremtræder som Darwins Forløber\footnote{Platons Stat. 5te Bog. (Heises Overs. af Staten. 2den Deel, S. 24--27).}. Den medfødte
Mulighed fører ikke direkte til Maalet; den forudsætter
Individets egen Stræben og virkeliggjøres ikke, naar Individet
ikke vil tage Sagen i sin Haand. Vistnok virker Energiens
høiere Udvikling tilbage paa det organisk-psykiske
Grundlag, som gaaer i Arv; det sees tydelig i de ædle
Slægter. Men dog er det ikke den lykkeligt udviklede Hjerne;
det er kun Betingelserne for dens lykkelige Udvikling, der
kunne arves.

Maudsley anvender ikke den samme Betragtningsmaade
paa Slægten som paa Individet. Hvad dette angaaer, hævder
han det medfødte Grundlags Betydning i Modsætning til de
ydre Paavirkningers Indflydelse; men ved hele Slægten mener
han ikke at behøve et saadant oprindeligt Grundlag; der tilskriver
han de ydre Paavirkninger en overvældende Betydning.
Og dog vil man vanskelig kunne fastholde en saadan modsat
Betragtningsmaade; man maa nødvendigviis anlægge en
beslægtet Maalestok begge Steder, saasandt som Slægtens
Væsen ikke kan være ueensartet med Individets. Den tydske
Materialist Büchner er her konsekventere end den engelske
Positivist; han vil end ikke indrømme et eiendommeligt, medfødt
Grundlag for Bevidsthedslivet, af Angst for, at de
„medfødte Ideer“ skulle indsmugles under denne Form. Mener
Maudsley af Erfaringsgrunde at maatte holde paa det Medfødtes
Betydning for Individets Vedkommende, maa han
ogsaa gjøre det for Slægtens Vedkommende.

Men det Medfødte kan i ethvert Tilfælde ikke siges uden
Videre at være Et med Individets Skjæbne, saaledes at
denne er afgjort ved Fødselen. Skjæbnen i sand Forstand,
den ethiske Skjæbne, Resultatet af Vexelvirkningen mellem
Individets bevidste Stræben og de indre og ydre Betingelser,
kan ikke være medfødt. Den bevidste Stræben er, selv om
% --- p. 331 ---
\opage{331}den maa bukke under for de ugunstige Betingelser som for
Naturnødvendighedens Magt, en ethisk Realitet. Kun den reent
idiotiske Existens er umiddelbart bestemt ved det medfødte
Grundlag; men denne Existens mangler saa ogsaa det Menneskeliges
% as printed (p. 331): „Karakteermærke“; not corrected
egenlige Karakteermærke. Maudsley erklærer med
Rette, at Idioten er berøvet sin Fødselsret som Menneske,
da han er født med en saadan Hjernemangel, at han slet
ikke eller dog kun ufuldkomment kan udføre nogen sjælelig
Funktion. Men hvor maa man forbauses ved at læse hans
følgende Ord: „Han (Idioten) selv har ingen Skyld i
denne Hjemsøgelse, eftersom han ikke har begaaet anden
Brøde end den, at han har sin Andeel i den menneskelige
Syndefuldhed.“ Hvad vil dette dog sige, naar det ikke er
en blot Levning af en theologisk Metafysik? At have Andeel
i Syndefuldheden og dog ikke at have begaaet nogen Brøde,
ja endog mangle Muligheden dertil, da man fra første Færd
er berøvet Fødselsretten som Menneske! Syndefuldhed kan
ligesaalidt arves, som den moralske Sands kan det. En
saadan Betragtningsmaade fører til en fuldstændig Sammenblanding
af medicinske og moralske Begreber. Modstykket
til Idiotens Syndefuldhed dannes af Opfattelsen af Forbrydelse
som Sindssygdom eller Epilepsi.

Og dog er den sidste Opfattelse rigtigere end den første;
den fordrer kun en nærmere Bestemmelse. Man maa her
ikke sammenblande det juridiske og det moralske Synspunkt.
Forbrydelsen kan godt være juridisk tilstede, uden at være
det moralsk. Man kan tænke sig en Existens, der om den
end ikke har mistet sin Fødselsret som Menneske, dog er
fremkommet og udviklet under Forhold, som med en Naturnødvendigheds
Magt have maattet trænge alt Høiere og
Ædlere tilbage, hvor oprindeligt Anlæg og ydre Vilkaar have
samvirket i den Grad, at den forbryderiske Handling næsten
kan siges at komme som et nødvendigt Produkt af dem, uden
at der har været nogen Mulighed for Individets egen Arbeiden,
da Bevidstheden slet ikke har kunnet vaagne. Et
saadant Individ kan være juridisk strafbart og dog moralsk
% --- p. 332 ---
\opage{332}undskyldeligt: juridisk gaaes der ud fra den Enkeltes Forhold
til Samfundet og dettes Ret overfor ham; moralsk rettes
Blikket ene paa Forholdet mellem den bevidste Stræben og
Betingelserne. Men den almindelige Opfattelse seer Sagen
saaledes, at en kriminel Forbrydelse er en høiere Grad af en
blot moralsk Forseelse; den seer ikke, at det egenlig ethiske
Moment ofte kan være aldeles underordnet i det juridisk
Strafbare. Her kan da Samfundet være i sin Ret ved at
fastholde den juridiske Betragtningsmaade; men den moralske
maa heller ikke oversees: den indskærper paa eengang en forsonende
Opfattelse af „Forbryderen“ og peger hen paa, hvor
den egenlige Forskjel mellem Godt og Ondt kommer frem.
Den medicinske Betragtningsmaade har ubestridelig Ret overfor
den juridiske, naar den hævder, „at en betydelig Deel af
Forbryderne ere aandssvage eller epileptiske eller stamme fra
Familier, i hvilke der er Sindssygdom, Epilepsi eller en
anden Nervesygdom“. Men den gjør sig ofte skyldig i den samme
Sammenblanding som den vulgære Betragtningsmaade, kun
paa anden Viis. Fordi der kan reises Tvivl om kriminelle
Forbrydelsers moralske Tilregnelighed i visse Tilfælde, derfor
er denne Tilregnelighed ikke ophævet aldeles. Som om
den kun fandt Anvendelse, hvor der juridisk er Tale om
Forbrydelse!

% p. 332: the mark after „praktiske,“ is a speck, not a hyphen or dash
Maudsley bestemmer Lægens Standpunkt som det praktiske,
hvor Sjæl og Legeme betragtes som en Væsens-Eenhed;
kun reent theoretisk kan man skille Sjælen fra Legemet og
fordybe sig i Selvbevidstheden. I praktisk Henseende blive
vi maaskee ogsaa mere enige med ham end i Theorien.
Det er udenfor al Tvivl, at Lægen er kaldet, og stedse
mere og mere vil blive kaldet til et praktisk Hverv af den
største humane Betydning. Men netop derfor er det ogsaa
af den største Betydning, at hans Dannelse og Synsmaade
ikke bliver eensidig. Ved det praktiske Standpunkt synes
Maudsley at forstaae det fysiologiske i Modsætning til, hvad
han kalder det metafysiske. Og dog er ogsaa det fysiologiske
Standpunkt endnu kun et theoretisk-eensidigt, naar det ikke
% --- p. 333 ---
\opage{333}forbindes med en virkelig, levende psykologisk Indsigt; det
har noksom viist sig i det Foregaaende. Vel erhverver
Lægen ofte en saadan psykologisk Indsigt, idet han ved den
skarpe Opmærksomhed for Legemets Tilstande ogsaa skærper
sit Blik for Iagttagelse af det Sjælelige; men indtrængende
og fuldstændig kan den dog først blive, naar det Psykologiske
erkjendes som en i sig eiendommelig Kreds af Love og Betingelser,
der ikke blot er et Aftryk af de legemlige Tilstande.
Men saavidt synes det endnu ikke at være kommet. „Det
er sjeldent,“ siger Stuart Mill, „at finde en Læge, der tillige
er Psykolog.“ Selv Positivismens Stifter, Auguste
Comte, deelte denne Mening; han betragtede Lægerne kun
som Veterinærer, fordi de ene bekymre sig om Menneskets
dyriske, ikke om hans menneskelige Natur.

Det i Sandhed praktiske Standpunkt viser ogsaa her
tilbage til den sande Theori, idet det fordrer alle Hensyn
respekterede. Derfor kunne Positivismen og Materialismen,
der kun have grebet een Side af Livet og Tilværelsens
Væsen og altsaa ere eensidige i Theorien, heller ikke føre til
den sande Praxis.

% p. 333: signature set right in bold Fraktur; below it a faint centred double rule (c. 2.2 cm, two lines c. 5 px apart), printing at about the density of the footnote rule on p. 330 (darkest row mean 163--167 against 157 there; the ghost text of p. 334 around it is much paler), so taken as the closing rule, as under „Constans.“ in the sibling article; the rest of the page carries only show-through of the following article, not transcribed
\begin{flushright}\textbf{Harald Høffding.}\end{flushright}
\begin{center}\rule{2.2cm}{0.4pt}\\[-10pt]\rule{2.2cm}{0.4pt}\end{center}

\end{document}
